\documentclass[a4paper,12pt]{article}

% ====================================================================
% SCHALTER: Lösungen ein-/ausblenden
% ====================================================================
\newif\ifloesungen
\loesungentrue
%\loesungenfalse

% ====================================================================
% SPRACHE / SCHRIFT
% ====================================================================
\usepackage[T1]{fontenc}
\usepackage[ngerman]{babel}
\usepackage{libertinus}

% ====================================================================
% TYPOGRAFIE & MATHE
% ====================================================================
\usepackage{amsmath}
\usepackage{microtype}
\usepackage{parskip}
\usepackage{enumitem}

\setlist[itemize]{
  topsep=0.2em,
  partopsep=0.1em,
  parsep=0.3em,
  itemsep=0.15em
}

\setlist[enumerate]{
  topsep=0.2em,
  partopsep=0.1em,
  parsep=0.3em,
  itemsep=0.15em
}

% ====================================================================
% LAYOUT
% ====================================================================
\usepackage[margin=2cm]{geometry}

% ====================================================================
% FARBEN
% ====================================================================
\usepackage[table]{xcolor}

\definecolor{PrimaryColor}{HTML}{226699}
\definecolor{SecondaryColor}{HTML}{55ACEE}
\definecolor{ExampleColor}{HTML}{78B159}
\definecolor{AlertColor}{HTML}{A0041E}
\definecolor{BlockBackground}{HTML}{EBF7FF}
\definecolor{PrimaryTitle}{HTML}{194C73}
\definecolor{DarkGrey}{HTML}{606060}
\definecolor{ExampleTitle}{HTML}{558C41} 
\definecolor{AlertTitle}{HTML}{A0041E}   
\definecolor{White}{HTML}{FFFFFF}
\definecolor{TextColor}{HTML}{000000}

\colorlet{VertexColor}{DarkGrey!20}

% ====================================================================
% HEADER / FOOTER
% ====================================================================
\usepackage{fancyhdr}
\pagestyle{fancy}
\setlength{\headheight}{14pt}

\fancyhf{}
\fancyhead[L]{\small Methoden und Theorien der KI}
\fancyhead[R]{\small INF2024}
\fancyfoot[R]{\thepage}
\renewcommand{\headrulewidth}{0.4pt}

\fancypagestyle{cover}{
  \fancyhf{}
  \renewcommand{\headrulewidth}{0pt}
}

% ====================================================================
% AUFGABEN / LÖSUNGEN
% ====================================================================
\usepackage{tcolorbox}
\tcbuselibrary{breakable}

\usepackage{chngcntr}

\newcounter{aufgabe}
\counterwithin{aufgabe}{section} % reset pro section
\renewcommand{\theaufgabe}{\thesection.\arabic{aufgabe}}

\newtcolorbox{aufgabebox}[1]{%
  breakable,
  colback=BlockBackground,
  colframe=PrimaryTitle,
  colbacktitle=PrimaryTitle,
  coltitle=White,
  fonttitle=\bfseries,
  title=Aufgabe~#1,
  boxrule=0.5pt,
  arc=0.7mm,
  toptitle=1mm,
  bottomtitle=1mm
}

\newcommand{\aufgabe}[1]{%
  \stepcounter{aufgabe}
  \begin{aufgabebox}{\theaufgabe}
    #1
  \end{aufgabebox}
}

\newcommand{\loesung}[1]{%
  \ifloesungen
    \medskip
    \noindent\textbf{Lösung: }\;#1
    \par\medskip
  \fi
}

% ====================================================================
% ALGORITHMEN
% ====================================================================
\usepackage[linesnumbered,ruled,vlined]{algorithm2e}
\setlength{\algomargin}{2em}
\SetAlgorithmName{Algorithmus}{Algorithmenverzeichnis}{List of Algorithms}
\renewcommand{\thealgocf}{}

\newcommand{\commentstyle}{\itshape\color{DarkGrey}}
\SetCommentSty{commentstyle}

% ====================================================================
% TIKZ / GRAPHEN
% ====================================================================
\usepackage{tikz}
\usetikzlibrary{
  automata,
  positioning,
  arrows.meta,
  trees,
  shapes,
  shapes.geometric,
  decorations.pathmorphing,
  decorations.pathreplacing,
  graphs,
  graphdrawing,
  intersections
}
\usegdlibrary{force}

% ====================================================================
% DOKUMENT
% ====================================================================
\begin{document}

% --------------------------------------------------------------------
% DECKBLATT
% --------------------------------------------------------------------
\thispagestyle{cover}

\vspace*{5cm}

\begin{center}
  {\Huge\bfseries Methoden und Theorien \\[0.2cm] der Künstlichen Intelligenz}\\[0.4cm]
  {\Large\itshape Übungen}\\[1.2cm]
  {\Large Studienjahrgang INF2024}\\[0.4cm]
  {\Large Autor: Markus Kaupp}
\end{center}

\newpage

% --------------------------------------------------------------------
% INHALT
% --------------------------------------------------------------------
  \section{Lineare Regression}
  \setcounter{aufgabe}{0}

  \aufgabe{
Gegeben seien folgende Testdatenpunkte aus $\mathbb{R} \times \mathbb{R}$:
\[
(0,\,4),\quad (1,\,5),\quad (2,\,7),\quad (3,\,12).
\]
Weiter sei die Regressionsfunktion $f: \mathbb{R}\to \mathbb{R}$ mit
\[
f(\mathbf{x}) = u_0 + \mathbf{b}\cdot\mathbf{x}, \qquad u_0 = 2,\ \mathbf{b} = (3)
\]
gegeben, also $f(x) = 2 + 3x$.

Bestimmen Sie die Summe der quadrierten Fehler $E$.
}

\loesung{
Es gilt mit $m$ Datenpunkten:
\[
E = \sum_{i=1}^{m}\bigl(y_i - f(\mathbf{x}_i)\bigr)^2.
\]

Berechnung der Vorhersagen und Fehler ($m=4$):
\begin{itemize}
  \item $\hat{y}=f(0) = 2$, \quad $y=4$, \quad $(y-\hat{y})^2 = 4$
  \item $\hat{y}=f(1) = 5$, \quad $y=5$, \quad $(y-\hat{y})^2 = 0$
  \item $\hat{y}=f(2) = 8$, \quad $y=7$, \quad $(y-\hat{y})^2 = 1$
  \item $\hat{y}=f(3) = 11$, \quad $y=12$, \quad $(y-\hat{y})^2 = 1$
\end{itemize}

\[
E = 4+0+1+1 = 6.
\]
}


\aufgabe{
Gegeben seien die Testdatenpunkte im Format $((x_1,x_2,x_3)^{\mathsf{T}},\,y)$ mit
Eingangsdimension $n=3$:
\[
((4,0,0)^{\mathsf{T}},\,8),\quad ((2,1,1)^{\mathsf{T}},\,5),\quad ((1,1,1)^{\mathsf{T}},\,4),\quad ((0,0,2)^{\mathsf{T}},\,3).
\]
Weiter sei die Regressionsfunktion $f: \mathbb{R}^3 \to \mathbb{R}$ mit
\[
f(\mathbf{x}) = u_0 + \mathbf{b}\cdot\mathbf{x}, \qquad
u_0 = 1,\ \mathbf{b} = (2,-1,1)^{\mathsf{T}}
\]
gegeben.

Bestimmen Sie die Summe der quadrierten Fehler $E$.
}

\loesung{
Es gilt mit $m$ Datenpunkten:
\[
E = \sum_{i=1}^{m}\bigl(y_i - f(\mathbf{x}_i)\bigr)^2.
\]

Berechnung der Vorhersagen und Fehler ($m=4$):
\begin{itemize}
  \item $\hat{y}=f(4,0,0) = 1+8-0+0 = 9$, \quad $y=8$, \quad $(y-\hat{y})^2 = 1$
  \item $\hat{y}=f(2,1,1) = 1+4-1+1 = 5$, \quad $y=5$, \quad $(y-\hat{y})^2 = 0$
  \item $\hat{y}=f(1,1,1) = 1+2-1+1 = 3$, \quad $y=4$, \quad $(y-\hat{y})^2 = 1$
  \item $\hat{y}=f(0,0,2) = 1+0-0+2 = 3$, \quad $y=3$, \quad $(y-\hat{y})^2 = 0$
\end{itemize}

\[
E = 1+0+1+0 = 2.
\]
}

\aufgabe{
Gegeben seien die folgenden Testdatenpaare $(\mathbf{x},y)$ mit
$\mathbf{x}\in\mathbb{R}^5$ und $y\in\mathbb{R}$:
\[
((1,0,2,1,0)^{\mathsf{T}},\,5),\quad
((2,1,0,0,1)^{\mathsf{T}},\,5),\quad
((0,2,1,1,1)^{\mathsf{T}},\,4),
\]
\[
((1,1,1,2,0)^{\mathsf{T}},\,7),\quad
((3,0,1,0,2)^{\mathsf{T}},\,0).
\]

Weiter sei die lineare Regressionsfunktion
$f:\mathbb{R}^5\to\mathbb{R}$ gegeben durch
\[
f(\mathbf{x}) = u_0 + \mathbf{b} \cdot \mathbf{x},
\]
wobei
\[
u_0 = 3 \qquad\text{und}\qquad \mathbf{b} =(1,2,-1,1,-2)^{\mathsf{T}}
\]

Bestimmen Sie die Summe der quadrierten Fehler.
}


\loesung{
Es gilt mit $m$ Datenpunkten:
\[
E = \sum_{i=1}^{m}\bigl(y_i - f(\mathbf{x}_i)\bigr)^2.
\]

Berechnung der Vorhersagen und Fehler ($m=5$):
\begin{itemize}
  \item $\hat{y}=f(1,0,2,1,0) = 3+1+0-2+1-0 = 3$, \quad $y=5$, \quad $(y-\hat{y})^2 = 4$
  \item $\hat{y}=f(2,1,0,0,1) = 3+2+2-0+0-2 = 5$, \quad $y=5$, \quad $(y-\hat{y})^2 = 0$
  \item $\hat{y}=f(0,2,1,1,1) = 3+0+4-1+1-2 = 5$, \quad $y=4$, \quad $(y-\hat{y})^2 = 1$
  \item $\hat{y}=f(1,1,1,2,0) = 3+1+2-1+2-0 = 7$, \quad $y=7$, \quad $(y-\hat{y})^2 = 0$
  \item $\hat{y}=f(3,0,1,0,2) = 3+3+0-1+0-4 = 1$, \quad $y=0$, \quad $(y-\hat{y})^2 = 1$
\end{itemize}

\[
E = 4+0+1+0+1 = 6.
\]
}


\aufgabe{
Gegeben seien die Datenpunkte
\[
(x_1,y_1)=(0,1),\qquad
(x_2,y_2)=(1,5),\qquad
(x_3,y_3)=(2,3).
\]

Gesucht ist die Regressionsgerade
\[
y = u_0 + u_1 x.
\]

\begin{enumerate}[label=\alph*)]
  \item Stellen Sie die Matrix $A$ und den Vektor $\mathbf{y}$ auf.
  \item Berechnen Sie $A^{\mathsf{T}} A$.
  \item Invertieren Sie $A^{\mathsf{T}} A$.
  \item Bestimmen Sie den Gewichtsvektor $\mathbf{u} = (A^{\mathsf{T}} A)^{-1} A^{\mathsf{T}} \mathbf{y}$.
\end{enumerate}
}

\loesung{
\begin{enumerate}[label=\alph*)]
  \item Jede Zeile enthält den Offset $1$ und den Merkmalswert $x$:
  \[
  A = \begin{pmatrix} 1 & 0 \\ 1 & 1 \\ 1 & 2 \end{pmatrix},
  \qquad
  \mathbf{y} = \begin{pmatrix} 1 \\ 5 \\ 3 \end{pmatrix}.
  \]

  \item 
  \[
  A^{\mathsf{T}} A =
  \begin{pmatrix} 1 & 1 & 1 \\ 0 & 1 & 2 \end{pmatrix}
  \begin{pmatrix} 1 & 0 \\ 1 & 1 \\ 1 & 2 \end{pmatrix}
  =
  \begin{pmatrix} 3 & 3 \\ 3 & 5 \end{pmatrix}.
  \]

  \item Determinante: $\det(A^{\mathsf{T}} A) = 3\cdot 5 - 3\cdot 3 = 6$.
  Für eine $2\times 2$-Matrix gilt
  $\begin{pmatrix} a & b \\ c & d \end{pmatrix}^{-1}
  = \frac{1}{ad-bc}\begin{pmatrix} d & -b \\ -c & a \end{pmatrix}$,
  also
  \[
  (A^{\mathsf{T}} A)^{-1}
  = \frac{1}{6}\begin{pmatrix} 5 & -3 \\ -3 & 3 \end{pmatrix}.
  \]

  \item 
  \[
  A^{\mathsf{T}} \mathbf{y}
  = \begin{pmatrix} 1 & 1 & 1 \\ 0 & 1 & 2 \end{pmatrix}
    \begin{pmatrix} 1 \\ 5 \\ 3 \end{pmatrix}
  = \begin{pmatrix} 1+5+3 \\ 0+5+6 \end{pmatrix}
  = \begin{pmatrix} 9 \\ 11 \end{pmatrix}.
  \]
  \[
  \mathbf{u}
  = \frac{1}{6}\begin{pmatrix} 5 & -3 \\ -3 & 3 \end{pmatrix}
    \begin{pmatrix} 9 \\ 11 \end{pmatrix}
  = \frac{1}{6}\begin{pmatrix} 45-33 \\ -27+33 \end{pmatrix}
  = \frac{1}{6}\begin{pmatrix} 12 \\ 6 \end{pmatrix}
  = \begin{pmatrix} 2 \\ 1 \end{pmatrix}.
  \]
  Das Modell lautet also $f(x) = 2 + x$.
\end{enumerate}
}


\aufgabe{
Gegeben seien folgende Beobachtungen aus $\mathbb{R}^2\times \mathbb{R}$ im Format $(x_1, x_2, y)$
\[
 \bigl\{(1,1,3),\,(2,1,6),\,(1,2,7),\,(2,2,6)\bigr\}.
\]
Gesucht ist die Regressionsebene 
\[f:\mathbb{R}^2 \to \mathbb{R}, \mathbf{x} \mapsto u_0+u_1x_1 + u_2 x_2 \] 
mit Gewichtsvektor $\mathbf{u} = (u_0, u_1, u_2)^{\mathsf{T}} \in \mathbb{R}^3$.\par\medskip

\begin{enumerate}[label=\alph*)]
  \item Stellen Sie die Matrix $A \in \mathbb{R}^{4\times 3}$ und den Zielvektor $\mathbf{y} \in \mathbb{R}^4$ auf, sodass gilt $\mathbf{y} = A\mathbf{u}$.
  \item Berechnen Sie $A^{\mathsf{T}} A$.
  \item Bestimmen Sie die Inverse $\bigl(A^{\mathsf{T}} A\bigr)^{-1}$.
  \item Berechnen Sie den Gewichtsvektor $\mathbf{u} = \bigl(A^{\mathsf{T}} A\bigr)^{-1} A^{\mathsf{T}} \mathbf{y}$.
\end{enumerate}
}

\loesung{
\begin{enumerate}[label=\alph*)]
  \item Jede Zeile enthält den Offset $1$ sowie die Merkmale $x_1, x_2$:
  \[
  A = \begin{pmatrix}
  1 & 1 & 1 \\
  1 & 2 & 1 \\
  1 & 1 & 2 \\
  1 & 2 & 2
  \end{pmatrix},
  \qquad
  \mathbf{y} = \begin{pmatrix} 3 \\ 6 \\ 7 \\ 6 \end{pmatrix}.
  \]

  \item Die Einträge ergeben sich aus den Spaltensummen und Kreuzsummen
  ($n=4$, $\sum x_1 = 6$, $\sum x_2 = 6$, $\sum x_1^2 = 10$, $\sum x_2^2 = 10$, $\sum x_1 x_2 = 9$):
  \[
  A^{\mathsf{T}} A =
  \begin{pmatrix}
  4 & 6 & 6 \\
  6 & 10 & 9 \\
  6 & 9 & 10
  \end{pmatrix}.
  \]

  \item Entwicklung nach der ersten Zeile:
  \[
  \det(A^{\mathsf{T}} A)
  = 4(10\cdot 10 - 9\cdot 9) - 6(6\cdot 10 - 9\cdot 6) + 6(6\cdot 9 - 10\cdot 6)
  = 4\cdot 19 - 6\cdot 6 + 6\cdot(-6) = 4.
  \]
  Mit der Adjunkten (Kofaktormatrix) folgt:
  \[
  (A^{\mathsf{T}} A)^{-1}
  = \frac{1}{4}
  \begin{pmatrix}
  19 & -6 & -6 \\
  -6 & 4 & 0 \\
  -6 & 0 & 4
  \end{pmatrix}.
  \]

  \item Zunächst
  \[
  A^{\mathsf{T}} \mathbf{y} =
  \begin{pmatrix}
  3+6+7+6 \\
  3+12+7+12 \\
  3+6+14+12
  \end{pmatrix}
  = \begin{pmatrix} 22 \\ 34 \\ 35 \end{pmatrix}.
  \]
  Damit
  \[
  \mathbf{u}
  = \frac{1}{4}
  \begin{pmatrix}
  19 & -6 & -6 \\
  -6 & 4 & 0 \\
  -6 & 0 & 4
  \end{pmatrix}
  \begin{pmatrix} 22 \\ 34 \\ 35 \end{pmatrix}
  = \frac{1}{4}\begin{pmatrix} 418-204-210 \\ -132+136 \\ -132+140 \end{pmatrix}
  = \frac{1}{4}\begin{pmatrix} 4 \\ 4 \\ 8 \end{pmatrix}
  = \begin{pmatrix} 1 \\ 1 \\ 2 \end{pmatrix}.
  \]
  Das Modell lautet also $f(x_1,x_2) = 1 + x_1 + 2x_2$.
\end{enumerate}
}






\section{Modellauswahl}
\setcounter{aufgabe}{0}

\aufgabe{
Gegeben seien die Testdatenpunkte im Format $((x_1,\dots,x_n),\,y)$ mit
Eingangsdimension $n=1$:
\[
((1),\,2),\quad ((2),\,5),\quad ((3),\,6),\quad ((4),\,9).
\]
Weiter sei die Regressionsfunktion
\[
f(\mathbf{x}) = u_0 + \mathbf{b}\cdot\mathbf{x}, \qquad u_0 = 1,\ \mathbf{b} = (2)
\]
gegeben, also $f(x) = 1 + 2x$.

Bestimmen Sie:
\begin{enumerate}[label=\alph*)]
  \item die Summe der quadrierten Abstände,
  \item den $R^2$-Score.
\end{enumerate}
}

\loesung{
Es bezeichne $m$ die Anzahl der Datenpunkte und $\hat{y}_i := f(\mathbf{x}_i)$ die
Vorhersage. Gemäß Definition des Bestimmtheitsmaßes gilt
\[
E_{\text{modell}} = \sum_{i=1}^{m}\bigl(y_i - \hat{y}_i\bigr)^2,
\qquad
R^2 = 1 - \frac{E_{\text{modell}}}{E_{\text{vergleich}}},
\qquad
E_{\text{vergleich}} = \sum_{i=1}^{m}\bigl(y_i - \bar{y}\bigr)^2,
\qquad
\bar{y} = \frac{1}{m}\sum_{i=1}^{m} y_i.
\]
$E_{\text{modell}}$ ist der Fehler des zu bewertenden Modells, $E_{\text{vergleich}}$
der Fehler des Vergleichsmodells $\hat{y}_i = \bar{y}$.

\medskip
Vorhersagen und Residuen ($m=4$):
\begin{itemize}
  \item $\hat{y}_1 = f(1) = 3$, \quad $y_1=2$, \quad $(y_1-\hat{y}_1)^2 = 1$
  \item $\hat{y}_2 = f(2) = 5$, \quad $y_2=5$, \quad $(y_2-\hat{y}_2)^2 = 0$
  \item $\hat{y}_3 = f(3) = 7$, \quad $y_3=6$, \quad $(y_3-\hat{y}_3)^2 = 1$
  \item $\hat{y}_4 = f(4) = 9$, \quad $y_4=9$, \quad $(y_4-\hat{y}_4)^2 = 0$
\end{itemize}

\begin{enumerate}[label=\alph*)]
  \item Summe der quadrierten Abstände:
  \[
  E_{\text{modell}} = 1+0+1+0 = 2.
  \]

  \item Mittelwert: $\bar{y} = \dfrac{2+5+6+9}{4} = \dfrac{22}{4} = 5{,}5$.

  Fehler des Vergleichsmodells:
  \[
  E_{\text{vergleich}}
  = (2-5{,}5)^2+(5-5{,}5)^2+(6-5{,}5)^2+(9-5{,}5)^2
  = 12{,}25+0{,}25+0{,}25+12{,}25 = 25.
  \]
  \[
  R^2 = 1 - \frac{2}{25} = \frac{23}{25} = 0{,}92.
  \]
\end{enumerate}
}


\aufgabe{
Gegeben seien die Testdatenpunkte im Format $((x_1,x_2,x_3),\,y)$ mit
Eingangsdimension $n=3$:
\[
((5,1,0),\,7),\quad ((1,4,2),\,8),\quad ((0,2,2),\,2),\quad ((4,0,2),\,3).
\]
Weiter sei die Regressionsfunktion
\[
f(\mathbf{x}) = u_0 + \mathbf{b}\cdot\mathbf{x}, \qquad
u_0 = 1,\ \mathbf{b} = (1,2,-1)
\]
gegeben, also $f(\mathbf{x}) = 1 + x_1 + 2x_2 - x_3$.

Bestimmen Sie:
\begin{enumerate}[label=\alph*)]
  \item die Summe der quadrierten Abstände,
  \item den $R^2$-Score.
\end{enumerate}
}

\loesung{
Mit $\hat{y}_i := f(\mathbf{x}_i)$ und $m$ Datenpunkten gilt
\[
E_{\text{modell}} = \sum_{i=1}^{m}\bigl(y_i - \hat{y}_i\bigr)^2,
\qquad
R^2 = 1 - \frac{E_{\text{modell}}}{E_{\text{vergleich}}},
\qquad
E_{\text{vergleich}} = \sum_{i=1}^{m}\bigl(y_i - \bar{y}\bigr)^2,
\qquad
\bar{y} = \frac{1}{m}\sum_{i=1}^{m} y_i.
\]

\medskip
Vorhersagen und Residuen ($m=4$):
\begin{itemize}
  \item $\hat{y}_1 = f(5,1,0) = 1+5+2-0 = 8$, \quad $y_1=7$, \quad $(y_1-\hat{y}_1)^2 = 1$
  \item $\hat{y}_2 = f(1,4,2) = 1+1+8-2 = 8$, \quad $y_2=8$, \quad $(y_2-\hat{y}_2)^2 = 0$
  \item $\hat{y}_3 = f(0,2,2) = 1+0+4-2 = 3$, \quad $y_3=2$, \quad $(y_3-\hat{y}_3)^2 = 1$
  \item $\hat{y}_4 = f(4,0,2) = 1+4+0-2 = 3$, \quad $y_4=3$, \quad $(y_4-\hat{y}_4)^2 = 0$
\end{itemize}

\begin{enumerate}[label=\alph*)]
  \item Summe der quadrierten Abstände:
  \[
  E_{\text{modell}} = 1+0+1+0 = 2.
  \]

  \item Mittelwert: $\bar{y} = \dfrac{7+8+2+3}{4} = \dfrac{20}{4} = 5$.

  Fehler des Vergleichsmodells:
  \[
  E_{\text{vergleich}}
  = (7-5)^2+(8-5)^2+(2-5)^2+(3-5)^2
  = 4+9+9+4 = 26.
  \]
  \[
  R^2 = 1 - \frac{2}{26} = \frac{12}{13} \approx 0{,}923.
  \]
\end{enumerate}
}


\aufgabe{
Gegeben seien die Testdatenpunkte im Format $((x_1,\dots,x_5),\,y)$ mit
Eingangsdimension $n=5$:
\[
((2,1,0,1,3),\,2),\quad
((1,0,2,1,1),\,7),\quad
((0,2,1,0,0),\,2),
\]
\[
((3,1,1,2,1),\,7),\quad
((1,1,0,0,1),\,2).
\]
Weiter sei die Regressionsfunktion
\[
f(\mathbf{x}) = u_0 + \mathbf{b}\cdot\mathbf{x}, \qquad
u_0 = 2,\ \mathbf{b} = (1,-1,2,1,-1)
\]
gegeben, also $f(\mathbf{x}) = 2 + x_1 - x_2 + 2x_3 + x_4 - x_5$.

Bestimmen Sie:
\begin{enumerate}[label=\alph*)]
  \item die Summe der quadrierten Abstände,
  \item den $R^2$-Score.
\end{enumerate}
}

\loesung{
Mit $\hat{y}_i := f(\mathbf{x}_i)$ und $m$ Datenpunkten gilt
\[
E_{\text{modell}} = \sum_{i=1}^{m}\bigl(y_i - \hat{y}_i\bigr)^2,
\qquad
R^2 = 1 - \frac{E_{\text{modell}}}{E_{\text{vergleich}}},
\qquad
E_{\text{vergleich}} = \sum_{i=1}^{m}\bigl(y_i - \bar{y}\bigr)^2,
\qquad
\bar{y} = \frac{1}{m}\sum_{i=1}^{m} y_i.
\]

\medskip
Vorhersagen und Residuen ($m=5$):
\begin{itemize}
  \item $\hat{y}_1 = f(2,1,0,1,3) = 2+2-1+0+1-3 = 1$, \quad $y_1=2$, \quad $(y_1-\hat{y}_1)^2 = 1$
  \item $\hat{y}_2 = f(1,0,2,1,1) = 2+1-0+4+1-1 = 7$, \quad $y_2=7$, \quad $(y_2-\hat{y}_2)^2 = 0$
  \item $\hat{y}_3 = f(0,2,1,0,0) = 2+0-2+2+0-0 = 2$, \quad $y_3=2$, \quad $(y_3-\hat{y}_3)^2 = 0$
  \item $\hat{y}_4 = f(3,1,1,2,1) = 2+3-1+2+2-1 = 7$, \quad $y_4=7$, \quad $(y_4-\hat{y}_4)^2 = 0$
  \item $\hat{y}_5 = f(1,1,0,0,1) = 2+1-1+0+0-1 = 1$, \quad $y_5=2$, \quad $(y_5-\hat{y}_5)^2 = 1$
\end{itemize}

\begin{enumerate}[label=\alph*)]
  \item Summe der quadrierten Abstände:
  \[
  E_{\text{modell}} = 1+0+0+0+1 = 2.
  \]

  \item Mittelwert: $\bar{y} = \dfrac{2+7+2+7+2}{5} = \dfrac{20}{5} = 4$.

  Fehler des Vergleichsmodells:
  \[
  E_{\text{vergleich}}
  = (2-4)^2+(7-4)^2+(2-4)^2+(7-4)^2+(2-4)^2
  = 4+9+4+9+4 = 30.
  \]
  \[
  R^2 = 1 - \frac{2}{30} = \frac{14}{15} \approx 0{,}933.
  \]
\end{enumerate}
}

\section{Entscheidungsbäume}
\setcounter{aufgabe}{0}


% ============================================================
% AUFGABE 1 (einfach)
% ============================================================
\aufgabe{
Gegeben sei der folgende Trainingsdatensatz. Zielattribut ist \textbf{Grillen}.
\begin{center}
{\scriptsize
\renewcommand{\arraystretch}{1.25}
\setlength{\tabcolsep}{6pt}
\rowcolors{2}{SecondaryColor!12}{white}
\begin{tabular}{c c c c}
\rowcolor{PrimaryColor}
\textcolor{white}{\textbf{Nr.}} &
\textcolor{white}{\textbf{Wochenende}} &
\textcolor{white}{\textbf{Regen}} &
\textcolor{white}{\textbf{Grillen}}\\
1 & Ja   & Nein & Ja\\
2 & Ja   & Nein & Ja\\
3 & Ja   & Ja   & Nein\\
4 & Nein & Nein & Nein\\
5 & Nein & Nein & Nein\\
6 & Nein & Ja   & Nein\\
\end{tabular}}
\end{center}

Erstellen Sie mit dem ID3-Algorithmus einen Entscheidungsbaum:
\begin{enumerate}[label=\alph*)]
  \item Berechnen Sie die bedingte Entropie $H(V \mid W)$ für beide Attribute
        und bestimmen Sie das Wurzelattribut.
  \item Setzen Sie das Verfahren rekursiv fort und zeichnen Sie den vollständigen
        Entscheidungsbaum.
\end{enumerate}
}

\loesung{
Es genügt, die bedingte Entropie $H(V \mid W)$ zu berechnen und das
\textbf{kleinste} zu wählen, da $H(V)$ am Knoten konstant ist.

\begin{enumerate}[label=\alph*)]
  \item Von $6$ Beispielen sind $2\times$ Grillen ja und $4\times$ nein.
  \[
    H(\text{Grillen} \mid \text{Wochenende}) = -\Bigl[\,
      \underbrace{\tfrac{3}{6}}_{\text{Wo. ja}}
      \bigl(\tfrac{2}{3}\log_2\tfrac{2}{3} + \tfrac{1}{3}\log_2\tfrac{1}{3}\bigr)
      +
      \underbrace{\tfrac{3}{6}}_{\text{Wo. nein}}
      \bigl(\tfrac{3}{3}\log_2\tfrac{3}{3}\bigr)
    \,\Bigr] \approx 0{,}4592
  \]
  \[
    H(\text{Grillen} \mid \text{Regen}) = -\Bigl[\,
      \underbrace{\tfrac{4}{6}}_{\text{R. nein}}
      \bigl(\tfrac{2}{4}\log_2\tfrac{2}{4} + \tfrac{2}{4}\log_2\tfrac{2}{4}\bigr)
      +
      \underbrace{\tfrac{2}{6}}_{\text{R. ja}}
      \bigl(\tfrac{2}{2}\log_2\tfrac{2}{2}\bigr)
    \,\Bigr] \approx 0{,}6667
  \]
  Kleinste bedingte Entropie: \textbf{Wochenende} $\Rightarrow$ Wurzelattribut.

  \item Der Zweig \textbf{Wochenende = nein} ist rein (nur Grillen nein), also ein
  Blatt. Im Zweig \textbf{Wochenende = ja} (Beispiele 1--3) verbleibt nur das
  Attribut \textbf{Regen}, das die Restmenge exakt trennt:
  \[
    \text{Regen nein} \Rightarrow \text{Grillen},
    \qquad
    \text{Regen ja} \Rightarrow \text{nicht Grillen}.
  \]

  \vspace{0.4em}
  \centering
  \resizebox{0.6\linewidth}{!}{%
  \begin{tikzpicture}[
    test/.style={draw=PrimaryColor, very thick, rounded corners=4pt,
                 align=center, inner sep=5pt, text width=2.0cm, font=\footnotesize},
    rad/.style={draw=ExampleTitle, very thick, ellipse, align=center,
                inner sep=2pt, text width=1.3cm, minimum height=1.1cm, font=\scriptsize},
    nicht/.style={draw=AlertTitle, very thick, ellipse, align=center,
                  inner sep=2pt, text width=1.3cm, minimum height=1.1cm, font=\scriptsize},
    kante/.style={->, >=Stealth, thick},
    lab/.style={font=\footnotesize\itshape, fill=white, inner sep=1pt}
  ]
    \node[test]  (wo)   at (0,0)       {Wochenende};
    \node[test]  (re)   at (-2.6,-1.8) {Regen};
    \node[nicht] (n1)   at (2.6,-1.8)  {nicht Grillen};
    \node[rad]   (g1)   at (-4.2,-3.6) {Grillen};
    \node[nicht] (n2)   at (-1.0,-3.6) {nicht Grillen};

    \draw[kante] (wo) -- (re) node[lab, midway] {ja};
    \draw[kante] (wo) -- (n1) node[lab, midway] {nein};
    \draw[kante] (re) -- (g1) node[lab, midway] {nein};
    \draw[kante] (re) -- (n2) node[lab, midway] {ja};
  \end{tikzpicture}}
\end{enumerate}
}


% ============================================================
% AUFGABE 2 (schwierig, binäre Attribute)
% ============================================================
\aufgabe{
Gegeben sei der folgende Trainingsdatensatz. Zielattribut ist \textbf{Joggen}.
Alle Attribute sind binär.
\begin{center}
{\scriptsize
\renewcommand{\arraystretch}{1.25}
\setlength{\tabcolsep}{6pt}
\rowcolors{2}{SecondaryColor!12}{white}
\begin{tabular}{c c c c c}
\rowcolor{PrimaryColor}
\textcolor{white}{\textbf{Nr.}} &
\textcolor{white}{\textbf{Regen}} &
\textcolor{white}{\textbf{Windig}} &
\textcolor{white}{\textbf{Kalt}} &
\textcolor{white}{\textbf{Joggen}}\\
1 & Ja   & Nein & Nein & Nein\\
2 & Ja   & Ja   & Ja   & Nein\\
3 & Ja   & Nein & Ja   & Nein\\
4 & Nein & Nein & Nein & Ja\\
5 & Nein & Nein & Ja   & Ja\\
6 & Nein & Nein & Ja   & Ja\\
7 & Nein & Ja   & Nein & Ja\\
8 & Nein & Ja   & Ja   & Nein\\
9 & Nein & Ja   & Ja   & Nein\\
\end{tabular}}
\end{center}

Erstellen Sie mit dem ID3-Algorithmus einen Entscheidungsbaum:
\begin{enumerate}[label=\alph*)]
  \item Bestimmen Sie über die bedingte Entropie $H(V \mid W)$ das Wurzelattribut.
  \item Führen Sie die Attributwahl rekursiv fort und zeichnen Sie den
        vollständigen Entscheidungsbaum.
\end{enumerate}
}

\loesung{
Von $9$ Beispielen sind $4\times$ Joggen ja und $5\times$ nein. Es genügt, die
bedingte Entropie $H(V \mid W)$ zu berechnen und das \textbf{kleinste} zu wählen.

\begin{enumerate}[label=\alph*)]
  \item \textbf{Wurzel:} bedingte Entropie je Attribut.
  \[
    H(V \mid \text{Regen}) = -\Bigl[\,
      \underbrace{\tfrac{3}{9}}_{\text{ja}}\bigl(\tfrac{3}{3}\log_2\tfrac{3}{3}\bigr)
      +
      \underbrace{\tfrac{6}{9}}_{\text{nein}}
      \bigl(\tfrac{4}{6}\log_2\tfrac{4}{6} + \tfrac{2}{6}\log_2\tfrac{2}{6}\bigr)
    \,\Bigr] \approx 0{,}6122
  \]
  \[
    H(V \mid \text{Windig}) = -\Bigl[\,
      \underbrace{\tfrac{5}{9}}_{\text{nein}}
      \bigl(\tfrac{3}{5}\log_2\tfrac{3}{5} + \tfrac{2}{5}\log_2\tfrac{2}{5}\bigr)
      +
      \underbrace{\tfrac{4}{9}}_{\text{ja}}
      \bigl(\tfrac{1}{4}\log_2\tfrac{1}{4} + \tfrac{3}{4}\log_2\tfrac{3}{4}\bigr)
    \,\Bigr] \approx 0{,}9000
  \]
  \[
    H(V \mid \text{Kalt}) = -\Bigl[\,
      \underbrace{\tfrac{3}{9}}_{\text{nein}}
      \bigl(\tfrac{2}{3}\log_2\tfrac{2}{3} + \tfrac{1}{3}\log_2\tfrac{1}{3}\bigr)
      +
      \underbrace{\tfrac{6}{9}}_{\text{ja}}
      \bigl(\tfrac{2}{6}\log_2\tfrac{2}{6} + \tfrac{4}{6}\log_2\tfrac{4}{6}\bigr)
    \,\Bigr] \approx 0{,}9183
  \]
  Kleinste bedingte Entropie: \textbf{Regen} $\Rightarrow$ Wurzelattribut.

  \item \textbf{Rekursion.} Der Zweig \textbf{Regen = ja} (Beispiele 1--3) ist
  rein (nur Joggen nein), also ein Blatt.

  \medskip
  Im Zweig \textbf{Regen = nein} (Beispiele 4--9; $4\times$ ja, $2\times$ nein)
  verbleiben \textbf{Windig} und \textbf{Kalt}:
  \[
    H(V \mid \text{Windig}) = -\Bigl[\,
      \underbrace{\tfrac{3}{6}}_{\text{nein}}\bigl(\tfrac{3}{3}\log_2\tfrac{3}{3}\bigr)
      +
      \underbrace{\tfrac{3}{6}}_{\text{ja}}
      \bigl(\tfrac{1}{3}\log_2\tfrac{1}{3} + \tfrac{2}{3}\log_2\tfrac{2}{3}\bigr)
    \,\Bigr] \approx 0{,}4591
  \]
  \[
    H(V \mid \text{Kalt}) = -\Bigl[\,
      \underbrace{\tfrac{2}{6}}_{\text{nein}}\bigl(\tfrac{2}{2}\log_2\tfrac{2}{2}\bigr)
      +
      \underbrace{\tfrac{4}{6}}_{\text{ja}}
      \bigl(\tfrac{2}{4}\log_2\tfrac{2}{4} + \tfrac{2}{4}\log_2\tfrac{2}{4}\bigr)
    \,\Bigr] \approx 0{,}6667
  \]
  Kleinste bedingte Entropie: \textbf{Windig}. Der Zweig \textbf{Windig = nein}
  ist rein (Joggen). Im Zweig \textbf{Windig = ja} (Beispiele 7--9) trennt das
  letzte Attribut \textbf{Kalt} exakt: nein $\Rightarrow$ Joggen,
  ja $\Rightarrow$ nicht Joggen.

  \vspace{0.4em}
  \centering
  \resizebox{0.85\linewidth}{!}{%
  \begin{tikzpicture}[
    test/.style={draw=PrimaryColor, very thick, rounded corners=4pt,
                 align=center, inner sep=5pt, text width=1.7cm, font=\footnotesize},
    rad/.style={draw=ExampleTitle, very thick, ellipse, align=center,
                inner sep=2pt, text width=1.3cm, minimum height=1.1cm, font=\scriptsize},
    nicht/.style={draw=AlertTitle, very thick, ellipse, align=center,
                  inner sep=2pt, text width=1.3cm, minimum height=1.1cm, font=\scriptsize},
    kante/.style={->, >=Stealth, thick},
    lab/.style={font=\footnotesize\itshape, fill=white, inner sep=1pt}
  ]
    \node[test]  (regen) at (0,0)       {Regen};
    \node[test]  (wind)  at (-3.2,-1.9) {Windig};
    \node[nicht] (n1)    at (3.2,-1.9)  {nicht Joggen};
    \node[rad]   (j1)    at (-5.4,-3.8) {Joggen};
    \node[test]  (kalt)  at (-1.6,-3.8) {Kalt};
    \node[rad]   (j2)    at (-3.2,-5.7) {Joggen};
    \node[nicht] (n2)    at (0.0,-5.7)  {nicht Joggen};

    \draw[kante] (regen) -- (wind) node[lab, midway] {nein};
    \draw[kante] (regen) -- (n1)   node[lab, midway] {ja};
    \draw[kante] (wind)  -- (j1)   node[lab, midway] {nein};
    \draw[kante] (wind)  -- (kalt) node[lab, midway] {ja};
    \draw[kante] (kalt)  -- (j2)   node[lab, midway] {nein};
    \draw[kante] (kalt)  -- (n2)   node[lab, midway] {ja};
  \end{tikzpicture}}
\end{enumerate}
}
\section{Naiver Bayes}
\setcounter{aufgabe}{0}

% ============================================================
% AUFGABE 1: Satz von Bayes + totale Wahrscheinlichkeit
% ============================================================
\aufgabe{
Ein Betrieb sortiert \textbf{Lachs} und \textbf{Barsch}. In der Produktion sind
$70\,\%$ der Fische Lachse und $30\,\%$ Barsche. Ein Helligkeitssensor meldet
\textbf{hell}: Lachse sind zu $80\,\%$ hell, Barsche nur zu $20\,\%$.

\begin{enumerate}[label=\alph*)]
  \item Berechnen Sie über den Satz der totalen Wahrscheinlichkeit die
        Wahrscheinlichkeit $P(\text{hell})$.
  \item Ein Fisch wird als \textbf{hell} gemessen. Bestimmen Sie mit dem Satz
        von Bayes $P(\text{Lachs} \mid \text{hell})$.
  \item Wie wahrscheinlich ist es, dass dieser Fisch ein \textbf{Barsch} ist?
\end{enumerate}
}

\loesung{
\textbf{Gegebene Größen}
\[
\begin{aligned}
P(\text{Lachs}) &= 0{,}70 & P(\text{Barsch}) &= 0{,}30\\
P(\text{hell} \mid \text{Lachs}) &= 0{,}80 &
P(\text{hell} \mid \text{Barsch}) &= 0{,}20
\end{aligned}
\]

\begin{enumerate}[label=\alph*)]
  \item Totale Wahrscheinlichkeit für \textbf{hell}:
  \[
    P(\text{hell}) = P(\text{hell} \mid \text{Lachs})\,P(\text{Lachs})
                   + P(\text{hell} \mid \text{Barsch})\,P(\text{Barsch})
  \]
  \[
    P(\text{hell}) = 0{,}80 \cdot 0{,}70 + 0{,}20 \cdot 0{,}30
                   = 0{,}56 + 0{,}06 = 0{,}62.
  \]

  \item Satz von Bayes:
  \[
    P(\text{Lachs} \mid \text{hell})
      = \frac{P(\text{hell} \mid \text{Lachs})\,P(\text{Lachs})}{P(\text{hell})}
      = \frac{0{,}56}{0{,}62} \approx 0{,}903.
  \]

  \item Gegenwahrscheinlichkeit:
  \[
    P(\text{Barsch} \mid \text{hell}) = 1 - P(\text{Lachs} \mid \text{hell})
      = \frac{0{,}06}{0{,}62} \approx 0{,}097.
  \]
\end{enumerate}
}


% ============================================================
% AUFGABE 2: Naive Bayes mit diskreten Merkmalen
% ============================================================
\aufgabe{
Gegeben sei folgender Trainingsdatensatz mit den binären Merkmalen
\textbf{Länge} und \textbf{Helligkeit} sowie der Klasse (Lachs/Barsch).
\begin{center}
{\scriptsize
\renewcommand{\arraystretch}{1.25}
\setlength{\tabcolsep}{6pt}
\rowcolors{2}{SecondaryColor!12}{white}
\begin{tabular}{c c c c}
\rowcolor{PrimaryColor}
\textcolor{white}{\textbf{Nr.}} &
\textcolor{white}{\textbf{Länge}} &
\textcolor{white}{\textbf{Helligkeit}} &
\textcolor{white}{\textbf{Klasse}}\\
1 & lang & hell   & Lachs\\
2 & lang & hell   & Lachs\\
3 & lang & dunkel & Lachs\\
4 & lang & hell   & Lachs\\
5 & kurz & hell   & Lachs\\
6 & lang & hell   & Lachs\\
7 & kurz & dunkel & Barsch\\
8 & kurz & dunkel & Barsch\\
9 & lang & dunkel & Barsch\\
10& kurz & hell   & Barsch\\
\end{tabular}}
\end{center}

Ein neuer Fisch $B$ hat \textbf{Länge $=$ lang} und \textbf{Helligkeit $=$ hell}.
Klassifizieren Sie $B$ mit dem Naive-Bayes-Klassifikator.
\begin{enumerate}[label=\alph*)]
  \item Schätzen Sie die A-priori-Wahrscheinlichkeiten und die bedingten
        Wahrscheinlichkeiten aus den Häufigkeiten.
  \item Berechnen Sie für beide Klassen den Term
        $P(K)\cdot P(\text{Länge} \mid K)\cdot P(\text{Helligkeit} \mid K)$
        und bestimmen Sie die Klasse von $B$.
\end{enumerate}
}

\loesung{
Für den Vergleich genügt der Zähler $P(K)\cdot\prod_i P(M_i \mid K)$, da der
Nenner $P(B)$ für beide Klassen gleich ist.

\begin{enumerate}[label=\alph*)]
  \item Von $10$ Fischen sind $6$ Lachse und $4$ Barsche:
  \[
    P(\text{Lachs}) = \tfrac{6}{10} = 0{,}6,
    \qquad
    P(\text{Barsch}) = \tfrac{4}{10} = 0{,}4.
  \]
  Bedingte Wahrscheinlichkeiten:
  \[
    P(\text{lang} \mid \text{Lachs}) = \tfrac{5}{6},
    \qquad
    P(\text{hell} \mid \text{Lachs}) = \tfrac{5}{6},
  \]
  \[
    P(\text{lang} \mid \text{Barsch}) = \tfrac{1}{4},
    \qquad
    P(\text{hell} \mid \text{Barsch}) = \tfrac{1}{4}.
  \]

  \item Zähler je Klasse:
  \[
    P(\text{Lachs})\cdot P(\text{lang}\mid\text{Lachs})\cdot P(\text{hell}\mid\text{Lachs})
      = 0{,}6 \cdot \tfrac{5}{6} \cdot \tfrac{5}{6}
      = 0{,}6 \cdot \tfrac{25}{36} \approx 0{,}4167.
  \]
  \[
    P(\text{Barsch})\cdot P(\text{lang}\mid\text{Barsch})\cdot P(\text{hell}\mid\text{Barsch})
      = 0{,}4 \cdot \tfrac{1}{4} \cdot \tfrac{1}{4}
      = 0{,}4 \cdot \tfrac{1}{16} = 0{,}025.
  \]
  Wegen $0{,}4167 > 0{,}025$ wird $B$ als \textbf{Lachs} klassifiziert.

  Normiert ergibt sich:
  \[
    P(\text{Lachs} \mid B) = \frac{0{,}4167}{0{,}4167 + 0{,}025}
      \approx 0{,}943,
    \qquad
    P(\text{Barsch} \mid B) \approx 0{,}057.
  \]
\end{enumerate}
}


% ============================================================
% AUFGABE 3: Naive Bayes mit Gauß-Dichte
% ============================================================
\aufgabe{
Für die Merkmale \textbf{Helligkeit} und \textbf{Länge} wurden pro Klasse
Normalverteilungen mit Mittelwert $\mu$ und Standardabweichung $\sigma$
geschätzt:
\begin{center}
{\scriptsize
\renewcommand{\arraystretch}{1.25}
\setlength{\tabcolsep}{6pt}
\rowcolors{2}{SecondaryColor!12}{white}
\begin{tabular}{c c c c c}
\rowcolor{PrimaryColor}
\textcolor{white}{\textbf{Klasse}} &
\textcolor{white}{\textbf{$\mu_{\text{Hell}}$}} &
\textcolor{white}{\textbf{$\sigma_{\text{Hell}}$}} &
\textcolor{white}{\textbf{$\mu_{\text{Länge}}$}} &
\textcolor{white}{\textbf{$\sigma_{\text{Länge}}$}}\\
Lachs  & 6 & 2 & 12 & 3\\
Barsch & 9 & 2 & 8  & 3\\
\end{tabular}}
\end{center}

Die Klassen sind gleich häufig, also $P(\text{Lachs}) = P(\text{Barsch}) = 0{,}5$.
Ein neuer Fisch $B$ hat \textbf{Helligkeit $= 7$} und \textbf{Länge $= 10$}.
Klassifizieren Sie $B$ mit dem Naive-Bayes-Klassifikator über die
Gauß-Dichtefunktion
\[
  f(x) = \frac{1}{\sigma\sqrt{2\pi}}\,
         \exp\!\Bigl(-\frac{(x-\mu)^2}{2\sigma^2}\Bigr).
\]
}

\loesung{
Wegen der Unabhängigkeitsannahme gilt
$P(K)\cdot f(\text{Hell}{=}7 \mid K)\cdot f(\text{Länge}{=}10 \mid K)$.

\medskip
\textbf{Dichtewerte Lachs}
\[
  f(\text{Hell}{=}7 \mid \text{Lachs})
    = \frac{1}{2\sqrt{2\pi}}\,\exp\!\Bigl(-\frac{(7-6)^2}{2\cdot 2^2}\Bigr)
    = 0{,}1995 \cdot e^{-0{,}125} \approx 0{,}1760,
\]
\[
  f(\text{Länge}{=}10 \mid \text{Lachs})
    = \frac{1}{3\sqrt{2\pi}}\,\exp\!\Bigl(-\frac{(10-12)^2}{2\cdot 3^2}\Bigr)
    = 0{,}1330 \cdot e^{-0{,}2222} \approx 0{,}1065.
\]

\textbf{Dichtewerte Barsch}
\[
  f(\text{Hell}{=}7 \mid \text{Barsch})
    = \frac{1}{2\sqrt{2\pi}}\,\exp\!\Bigl(-\frac{(7-9)^2}{2\cdot 2^2}\Bigr)
    = 0{,}1995 \cdot e^{-0{,}5} \approx 0{,}1210,
\]
\[
  f(\text{Länge}{=}10 \mid \text{Barsch})
    = \frac{1}{3\sqrt{2\pi}}\,\exp\!\Bigl(-\frac{(10-8)^2}{2\cdot 3^2}\Bigr)
    = 0{,}1330 \cdot e^{-0{,}2222} \approx 0{,}1065.
\]

\medskip
\textbf{Zähler je Klasse}
\[
  P(\text{Lachs})\cdot 0{,}1760 \cdot 0{,}1065
    = 0{,}5 \cdot 0{,}01874 \approx 0{,}00937,
\]
\[
  P(\text{Barsch})\cdot 0{,}1210 \cdot 0{,}1065
    = 0{,}5 \cdot 0{,}01288 \approx 0{,}00644.
\]
Wegen $0{,}00937 > 0{,}00644$ wird $B$ als \textbf{Lachs} klassifiziert.

Normiert:
\[
  P(\text{Lachs} \mid B) \approx \frac{0{,}00937}{0{,}00937 + 0{,}00644}
    \approx 0{,}593,
  \qquad
  P(\text{Barsch} \mid B) \approx 0{,}407.
\]
}
\section{SVM}
\setcounter{aufgabe}{0}

% ============================================================
% AUFGABE 1: Entscheidungsfunktion, Klassifikation, Abstand
% ============================================================
\aufgabe{
Gegeben sei im $\mathbb{R}^2$ eine Trennebene mit Normalenvektor
$\mathbf{w} = \begin{pmatrix} 3 \\ 4 \end{pmatrix}$ und $w_0 = -10$. Die
Entscheidungsfunktion lautet $g(\mathbf{x}) = w_0 + \mathbf{w}\cdot\mathbf{x}$
mit
\[
  \mathbf{w}\cdot\mathbf{x} + w_0 > 0 \Rightarrow \text{Klasse A},
  \qquad
  \mathbf{w}\cdot\mathbf{x} + w_0 < 0 \Rightarrow \text{Klasse B}.
\]

\begin{enumerate}[label=\alph*)]
  \item Klassifizieren Sie die Punkte
        $\mathbf{x}_1 = \begin{pmatrix} 2 \\ 2 \end{pmatrix}$ und
        $\mathbf{x}_2 = \begin{pmatrix} 1 \\ 1 \end{pmatrix}$.
  \item Berechnen Sie den geometrischen Abstand $r = \dfrac{g(\mathbf{x_1})}{\lVert\mathbf{w}\rVert}$
        von $\mathbf{x}_1$ zur Trennebene.
\end{enumerate}
}

\loesung{
\begin{enumerate}[label=\alph*)]
  \item Für $\mathbf{x}_1$:
  \[
    g(\mathbf{x}_1) = 3\cdot 2 + 4\cdot 2 - 10 = 6 + 8 - 10 = 4 > 0
    \;\Rightarrow\; \textbf{Klasse A}.
  \]
  Für $\mathbf{x}_2$:
  \[
    g(\mathbf{x}_2) = 3\cdot 1 + 4\cdot 1 - 10 = 3 + 4 - 10 = -3 < 0
    \;\Rightarrow\; \textbf{Klasse B}.
  \]

  \item Mit $\lVert\mathbf{w}\rVert = \sqrt{3^2 + 4^2} = \sqrt{25} = 5$:
  \[
    r = \frac{g(\mathbf{x}_1)}{\lVert\mathbf{w}\rVert} = \frac{4}{5} = 0{,}8.
  \]
\end{enumerate}
}


% ============================================================
% AUFGABE 2: Margin und Skalierung
% ============================================================
\aufgabe{
Eine SVM habe den Normalenvektor
$\mathbf{w} = \begin{pmatrix} 1 \\ -2 \\ 2 \end{pmatrix}$ und $w_0 = 3$.
\begin{enumerate}[label=\alph*)]
  \item Berechnen Sie den Margin $\rho = \dfrac{1}{\lVert\mathbf{w}\rVert}$.
  \item Ein Support Vector $\mathbf{x}_s$ der Klasse $y = +1$ erfüllt die
        Bedingung mit Gleichheit, also
        $y\cdot(\mathbf{w}\cdot\mathbf{x}_s + w_0) = 1$. Prüfen Sie, ob
        $\mathbf{x}_s = \begin{pmatrix} 0 \\ 0 \\ -1 \end{pmatrix}$ ein
        solcher Support Vector ist.
  \item Erläutern Sie, warum ein Verdoppeln von $\mathbf{w}$ und $w_0$
        dieselbe Trennebene, aber einen anderen Wert von $\lVert\mathbf{w}\rVert$
        liefert.
\end{enumerate}
}

\loesung{
\begin{enumerate}[label=\alph*)]
  \item $\lVert\mathbf{w}\rVert = \sqrt{1^2 + (-2)^2 + 2^2} = \sqrt{9} = 3$, also
  \[
    \rho = \frac{1}{\lVert\mathbf{w}\rVert} = \frac{1}{3} \approx 0{,}333.
  \]

  \item Einsetzen:
  \[
    \mathbf{w}\cdot\mathbf{x}_s + w_0 = 1\cdot 0 + (-2)\cdot 0 + 2\cdot(-1) + 3 = 1.
  \]
  Mit $y = +1$ gilt $y\cdot(\mathbf{w}\cdot\mathbf{x}_s + w_0) = 1$. Die
  Bedingung ist mit \textbf{Gleichheit} erfüllt, $\mathbf{x}_s$ ist ein
  \textbf{Support Vector}.

  \item Ein Vielfaches $\alpha\mathbf{w}$, $\alpha w_0$ liefert dieselbe Menge
  $\{\mathbf{x} : \alpha\mathbf{w}\cdot\mathbf{x} + \alpha w_0 = 0\}$, also
  \textbf{dieselbe Ebene}, da die Nullstellenmenge invariant unter Skalierung
  ist. Der Betrag $\lVert\alpha\mathbf{w}\rVert = |\alpha|\,\lVert\mathbf{w}\rVert$
  ändert sich jedoch. Deshalb wird die Skalierung über
  $\rho\cdot\lVert\mathbf{w}\rVert = 1$ fixiert.
\end{enumerate}
}


% ============================================================
% AUFGABE 3: Optimales w aus Lambda, duale Klassifikation
% ============================================================
\aufgabe{
Ein linear trennbares Problem im $\mathbb{R}^2$ habe drei Support Vectors:
\begin{center}
{\scriptsize
\renewcommand{\arraystretch}{1.25}
\setlength{\tabcolsep}{6pt}
\rowcolors{2}{SecondaryColor!12}{white}
\begin{tabular}{c c c c}
\rowcolor{PrimaryColor}
\textcolor{white}{\textbf{$i$}} &
\textcolor{white}{\textbf{$\mathbf{x}_i$}} &
\textcolor{white}{\textbf{$y_i$}} &
\textcolor{white}{\textbf{$\lambda_i$}}\\
1 & $(2, 0)^{\mathsf{T}}$ & $+1$ & $1$\\
2 & $(0, 2)^{\mathsf{T}}$ & $+1$ & $1$\\
3 & $(0, 0)^{\mathsf{T}}$ & $-1$ & $2$\\
\end{tabular}}
\end{center}

\begin{enumerate}[label=\alph*)]
  \item Prüfen Sie die duale Nebenbedingung $\sum_i \lambda_i y_i = 0$.
  \item Berechnen Sie $\mathbf{w}^{*} = \sum_i \lambda_i y_i \mathbf{x}_i$.
  \item Klassifizieren Sie das neue Sample
        $\mathbf{z} = \begin{pmatrix} 3 \\ 3 \end{pmatrix}$ über das Vorzeichen
        von $g(\mathbf{z}) = \mathbf{w}^{*}\cdot\mathbf{z} + w_0^{*}$, wobei
        $w_0^{*} = -1$ gegeben ist.
\end{enumerate}
}

\loesung{
\begin{enumerate}[label=\alph*)]
  \item Duale Nebenbedingung:
  \[
    \sum_i \lambda_i y_i = 1\cdot(+1) + 1\cdot(+1) + 2\cdot(-1) = 1 + 1 - 2 = 0. \checkmark
  \]

  \item Optimales $\mathbf{w}^{*}$:
  \[
    \mathbf{w}^{*} = 1\cdot(+1)\begin{pmatrix} 2 \\ 0 \end{pmatrix}
                   + 1\cdot(+1)\begin{pmatrix} 0 \\ 2 \end{pmatrix}
                   + 2\cdot(-1)\begin{pmatrix} 0 \\ 0 \end{pmatrix}
                   = \begin{pmatrix} 2 \\ 2 \end{pmatrix}.
  \]

  \item Entscheidungsfunktion:
  \[
    g(\mathbf{z}) = \mathbf{w}^{*}\cdot\mathbf{z} + w_0^{*}
                  = 2\cdot 3 + 2\cdot 3 - 1 = 11 > 0
    \;\Rightarrow\; \textbf{Klasse A} \;(y = +1).
  \]
\end{enumerate}
}

% ============================================================
% AUFGABE 4: Klassifikation mit polynomialem Kernel
% ============================================================
\aufgabe{
Eine SVM mit \textbf{polynomialem Kernel}
$K(\mathbf{x}_i, \mathbf{x}_j) = (\mathbf{x}_i\cdot\mathbf{x}_j + c)^d$
und $c = 1$, $d = 2$ klassifiziert ein neues Sample $\mathbf{z}$ über
\[
  g(\mathbf{z}) = \sum_{i} \lambda_i\, y_i\, K(\mathbf{x}_i, \mathbf{z}) + w_0^{*}.
\]
Gegeben sind die Support Vectors mit $w_0^{*} = -1$:
\begin{center}
{\scriptsize
\renewcommand{\arraystretch}{1.25}
\setlength{\tabcolsep}{6pt}
\rowcolors{2}{SecondaryColor!12}{white}
\begin{tabular}{c c c c}
\rowcolor{PrimaryColor}
\textcolor{white}{\textbf{$i$}} &
\textcolor{white}{\textbf{$\mathbf{x}_i$}} &
\textcolor{white}{\textbf{$y_i$}} &
\textcolor{white}{\textbf{$\lambda_i$}}\\
1 & $(1, 0)^{\mathsf{T}}$ & $+1$ & $0{,}5$\\
2 & $(0, 1)^{\mathsf{T}}$ & $+1$ & $0{,}5$\\
3 & $(1, 1)^{\mathsf{T}}$ & $-1$ & $1$\\
\end{tabular}}
\end{center}

Klassifizieren Sie das Sample $\mathbf{z} = \begin{pmatrix} 2 \\ 0 \end{pmatrix}$.
}

\loesung{
\textbf{Kernelwerte} $K(\mathbf{x}_i, \mathbf{z}) = (\mathbf{x}_i\cdot\mathbf{z} + 1)^2$:
\[
  \mathbf{x}_1\cdot\mathbf{z} = 1\cdot 2 + 0\cdot 0 = 2,
  \qquad K(\mathbf{x}_1, \mathbf{z}) = (2+1)^2 = 9.
\]
\[
  \mathbf{x}_2\cdot\mathbf{z} = 0\cdot 2 + 1\cdot 0 = 0,
  \qquad K(\mathbf{x}_2, \mathbf{z}) = (0+1)^2 = 1.
\]
\[
  \mathbf{x}_3\cdot\mathbf{z} = 1\cdot 2 + 1\cdot 0 = 2,
  \qquad K(\mathbf{x}_3, \mathbf{z}) = (2+1)^2 = 9.
\]

\textbf{Entscheidungsfunktion:}
\[
  g(\mathbf{z}) = 0{,}5\cdot(+1)\cdot 9 + 0{,}5\cdot(+1)\cdot 1 + 1\cdot(-1)\cdot 9 - 1
\]
\[
  g(\mathbf{z}) = 4{,}5 + 0{,}5 - 9 - 1 = -5 < 0
  \;\Rightarrow\; \textbf{Klasse B} \;(y = -1).
\]
}


% ============================================================
% AUFGABE 4: Klassifikation mit RBF-Kernel
% ============================================================
\aufgabe{
Eine SVM mit \textbf{RBF-Kernel}
$K(\mathbf{x}, \mathbf{y}) = e^{-\gamma\,\lVert\mathbf{x}-\mathbf{y}\rVert^2}$
und $\gamma = 0{,}5$ klassifiziert ein neues Sample $\mathbf{z}$ über
\[
  g(\mathbf{z}) = \sum_{i} \lambda_i\, y_i\, K(\mathbf{x}_i, \mathbf{z}) + w_0^{*}.
\]
Gegeben sind die Support Vectors mit $w_0^{*} = 0$:
\begin{center}
{\scriptsize
\renewcommand{\arraystretch}{1.25}
\setlength{\tabcolsep}{6pt}
\rowcolors{2}{SecondaryColor!12}{white}
\begin{tabular}{c c c c}
\rowcolor{PrimaryColor}
\textcolor{white}{\textbf{$i$}} &
\textcolor{white}{\textbf{$\mathbf{x}_i$}} &
\textcolor{white}{\textbf{$y_i$}} &
\textcolor{white}{\textbf{$\lambda_i$}}\\
1 & $(0, 0)^{\mathsf{T}}$ & $+1$ & $1$\\
2 & $(2, 2)^{\mathsf{T}}$ & $-1$ & $1$\\
\end{tabular}}
\end{center}

Klassifizieren Sie das Sample $\mathbf{z} = \begin{pmatrix} 1 \\ 0 \end{pmatrix}$.
}

\loesung{
\textbf{Quadratische Abstände} $\lVert\mathbf{x}_i - \mathbf{z}\rVert^2$:
\[
  \lVert\mathbf{x}_1 - \mathbf{z}\rVert^2 = (0-1)^2 + (0-0)^2 = 1,
\]
\[
  \lVert\mathbf{x}_2 - \mathbf{z}\rVert^2 = (2-1)^2 + (2-0)^2 = 1 + 4 = 5.
\]

\textbf{Kernelwerte} $K(\mathbf{x}_i, \mathbf{z}) = e^{-0{,}5\,\lVert\mathbf{x}_i - \mathbf{z}\rVert^2}$:
\[
  K(\mathbf{x}_1, \mathbf{z}) = e^{-0{,}5\cdot 1} = e^{-0{,}5} \approx 0{,}607,
\]
\[
  K(\mathbf{x}_2, \mathbf{z}) = e^{-0{,}5\cdot 5} = e^{-2{,}5} \approx 0{,}082.
\]

\textbf{Entscheidungsfunktion:}
\[
  g(\mathbf{z}) = 1\cdot(+1)\cdot 0{,}607 + 1\cdot(-1)\cdot 0{,}082 + 0
\]
\[
  g(\mathbf{z}) = 0{,}607 - 0{,}082 = 0{,}525 > 0
  \;\Rightarrow\; \textbf{Klasse A} \;(y = +1).
\]

\textbf{Deutung:} $\mathbf{z}$ liegt näher am Support Vector der Klasse A
(Abstand$^2 = 1$) als an dem der Klasse B (Abstand$^2 = 5$). Da der RBF-Kernel
Nähe als Ähnlichkeit misst, überwiegt der Beitrag der Klasse A.
}


% ============================================================
% AUFGABE 6: RBF-Kernel und Schlupfvariablen
% ============================================================
\aufgabe{
\begin{enumerate}[label=\alph*)]
  \item Der RBF-Kernel ist definiert durch
        $K(\mathbf{x}, \mathbf{y}) = e^{-\gamma\,\lVert\mathbf{x}-\mathbf{y}\rVert^2}$.
        Berechnen Sie mit $\gamma = 0{,}5$ die Ähnlichkeit
        $K(\mathbf{x}, \mathbf{y})$ für
        $\mathbf{x} = \begin{pmatrix} 1 \\ 2 \end{pmatrix}$,
        $\mathbf{y} = \begin{pmatrix} 3 \\ 2 \end{pmatrix}$.
  \item Geben Sie $K(\mathbf{x}, \mathbf{x})$ an und begründen Sie den Wert.
  \item Ein Punkt mit Label $y_i = +1$ erfüllt
        $y_i\cdot(w_0 + \mathbf{w}\cdot\mathbf{x}_i) = -0{,}3$. Bestimmen Sie die
        Schlupfvariable $\xi_i$ aus $y_i\cdot(w_0 + \mathbf{w}\cdot\mathbf{x}_i) \geq 1 - \xi_i$
        (mit Gleichheit) und entscheiden Sie, ob der Punkt korrekt klassifiziert
        ist.
\end{enumerate}
}

\loesung{
\begin{enumerate}[label=\alph*)]
  \item Quadratischer Abstand:
  \[
    \lVert\mathbf{x}-\mathbf{y}\rVert^2 = (1-3)^2 + (2-2)^2 = 4 + 0 = 4.
  \]
  \[
    K(\mathbf{x}, \mathbf{y}) = e^{-0{,}5\cdot 4} = e^{-2} \approx 0{,}135.
  \]

  \item Am selben Punkt ist der Abstand null:
  \[
    K(\mathbf{x}, \mathbf{x}) = e^{-\gamma\cdot 0} = e^{0} = 1.
  \]
  Der Kernel misst Ähnlichkeit; identische Punkte haben maximale Ähnlichkeit $1$.

  \item Aus $y_i\cdot(w_0 + \mathbf{w}\cdot\mathbf{x}_i) = 1 - \xi_i$ folgt
  \[
    \xi_i = 1 - \bigl(-0{,}3\bigr) = 1{,}3.
  \]
  Wegen $\xi_i \geq 1$ liegt eine \textbf{Fehlklassifikation} vor.
\end{enumerate}
}

\section{Neuronale Netze}
\setcounter{aufgabe}{0}
% ===================================================================
% Aufgaben zu Neuronalen Netzen
% Themen: Perzeptron, Hebbsche Lernregel, Boolesche Funktionen,
%         Aktivierungsfunktionen, MLP-Vorwaertsrechnung,
%         Gradientenabstieg, Backpropagation
% ===================================================================

% -------------------------------------------------------------------
% 1) Perzeptron: Klassifikation und Trenngerade
% -------------------------------------------------------------------
\aufgabe{
Gegeben sei ein einlagiges Perzeptron mit einem Ausgabeneuron und dem
Eingaberaum $\mathbb{R}^2$. Der Gewichtsvektor sei
$\mathbf{w} = \begin{pmatrix} 2 \\ -1 \end{pmatrix} \in \mathbb{R}^2$, der
Schwellenwert $\theta = 1$. Die Netzeingabe ist
\[
  h : \mathbb{R}^2 \to \mathbb{R}, \qquad
  h(\mathbf{x}) = \mathbf{w} \cdot \mathbf{x},
\]
die Ausgabe $y = g\bigl(h(\mathbf{x})\bigr)$ mit der Heaviside-Funktion
\[
  g : \mathbb{R} \to \{0,1\}, \qquad
  g(h) =
  \begin{cases}
    1, & h > \theta,\\
    0, & h \leq \theta.
  \end{cases}
\]

\begin{enumerate}[label=\alph*)]
  \item Bestimmen Sie die Ausgabe für
        $\mathbf{x}_1 = \begin{pmatrix} 1 \\ 0 \end{pmatrix}$,
        $\mathbf{x}_2 = \begin{pmatrix} 0 \\ 2 \end{pmatrix}$ und
        $\mathbf{x}_3 = \begin{pmatrix} 2 \\ 3 \end{pmatrix}$.
  \item Geben Sie die Trenngerade $h(\mathbf{x}) - \theta = 0$ in der Form
        $x_2 = a\,x_1 + b$ an.
  \item Berechnen Sie den Abstand des Punktes $\mathbf{x}_1$ zur Trenngeraden.
\end{enumerate}
}

\loesung{
\begin{enumerate}[label=\alph*)]
  \item Es gilt $h(\mathbf{x}) = 2x_1 - x_2$:
  \[
    h(\mathbf{x}_1) = 2 - 0 = 2 > 1 \;\Rightarrow\; y = \textbf{1}
  \]
  \[
    h(\mathbf{x}_2) = 0 - 2 = -2 \leq 1 \;\Rightarrow\; y = \textbf{0}
  \]
  \[
    h(\mathbf{x}_3) = 4 - 3 = 1 \leq 1 \;\Rightarrow\; y = \textbf{0}
  \]
  Der Punkt $\mathbf{x}_3$ liegt genau \textbf{auf} der Trenngeraden und wird
  wegen $h \leq \theta$ der Klasse $0$ zugeordnet.

  \item Aus $2x_1 - x_2 - 1 = 0$ folgt
  \[
    x_2 = 2x_1 - 1.
  \]

  \item Mit $\lVert\mathbf{w}\rVert = \sqrt{2^2 + (-1)^2} = \sqrt{5}$:
  \[
    r = \frac{h(\mathbf{x}_1) - \theta}{\lVert\mathbf{w}\rVert}
      = \frac{2 - 1}{\sqrt{5}} \approx 0{,}447.
  \]
\end{enumerate}
}

% -------------------------------------------------------------------
% 2) Hebbsche Lernregel
% -------------------------------------------------------------------
\aufgabe{
Ein Perzeptron mit Bias-Neuron verwendet den erweiterten Eingabevektor
$\mathbf{x} = (x_0, x_1, x_2)^{\mathsf{T}}$ mit $x_0 = 1$ und den
Gewichtsvektor
\[
  \mathbf{w} = (-0{,}2;\; 0{,}5;\; -0{,}3)^{\mathsf{T}},
  \qquad \eta = 0{,}2 .
\]
Die Ausgabe ist $\hat{y} = \text{heaviside}(\mathbf{w}\cdot\mathbf{x})$, d. h.
$\hat{y} = 1$ für $\mathbf{w}\cdot\mathbf{x} > 0$, sonst $\hat{y} = 0$.

Führen Sie zwei Schritte der Hebbschen Lernregel
$\Delta w_j = \eta\cdot\text{error}\cdot x_j$ durch:

\begin{enumerate}[label=\alph*)]
  \item Beispiel $\mathbf{x} = (1, 1, 1)^{\mathsf{T}}$ mit $y = 1$.
  \item Beispiel $\mathbf{x} = (1, 1, 0)^{\mathsf{T}}$ mit $y = 0$.
  \item Geben Sie die resultierende Trenngerade in der Form
        $x_2 = a\,x_1 + b$ an.
\end{enumerate}
}

\loesung{
\begin{enumerate}[label=\alph*)]
  \item Netzeingabe:
  \[
    \mathbf{w}\cdot\mathbf{x} = -0{,}2 + 0{,}5 - 0{,}3 = 0
    \;\Rightarrow\; \hat{y} = 0, \qquad \text{error} = 1 - 0 = 1 .
  \]
  Gewichtsänderungen:
  \[
    \Delta w_j = 0{,}2\cdot 1\cdot x_j = 0{,}2
    \quad\text{für } j = 0,1,2
  \]
  \[
    \mathbf{w} = (0;\; 0{,}7;\; -0{,}1)^{\mathsf{T}}
  \]

  \item Netzeingabe:
  \[
    \mathbf{w}\cdot\mathbf{x} = 0 + 0{,}7 + 0 = 0{,}7 > 0
    \;\Rightarrow\; \hat{y} = 1, \qquad \text{error} = 0 - 1 = -1 .
  \]
  Gewichtsänderungen:
  \[
    \Delta w_0 = -0{,}2,\quad \Delta w_1 = -0{,}2,\quad \Delta w_2 = 0
  \]
  \[
    \mathbf{w} = (-0{,}2;\; 0{,}5;\; -0{,}1)^{\mathsf{T}}
  \]

  \item Aus $-0{,}2 + 0{,}5x_1 - 0{,}1x_2 = 0$ folgt nach Multiplikation
  mit $10$:
  \[
    5x_1 - x_2 - 2 = 0
    \quad\Longleftrightarrow\quad
    x_2 = 5x_1 - 2 .
  \]
\end{enumerate}
}

% -------------------------------------------------------------------
% 3) Boolesche Funktionen mit einem Perzeptron
% -------------------------------------------------------------------
\aufgabe{
Gegeben sei ein Perzeptron mit $w_1 = w_2 = 1$ und der Ausgabe
$y = 1$ genau dann, wenn $w_1x_1 + w_2x_2 > \theta$.

\begin{enumerate}[label=\alph*)]
  \item Welche boolesche Funktion realisiert das Perzeptron für
        $\theta = 1{,}5$? Stellen Sie die vollständige Wertetabelle auf.
  \item Für welche Werte von $\theta$ realisiert dasselbe Netz die
        OR-Funktion?
  \item Zeigen Sie rechnerisch, dass XOR mit $\theta \geq 0$ nicht
        realisierbar ist.
\end{enumerate}
}

\loesung{
\begin{enumerate}[label=\alph*)]
  \item Mit $h = x_1 + x_2$:
  \[
    \begin{array}{c|c|c|c}
      x_1 & x_2 & h & y\\\hline
      0 & 0 & 0 & 0\\
      0 & 1 & 1 & 0\\
      1 & 0 & 1 & 0\\
      1 & 1 & 2 & 1
    \end{array}
  \]
  Nur $h = 2 > 1{,}5$ liefert $y = 1$, also die \textbf{AND}-Funktion.

  \item Gefordert ist $0 \leq \theta$ (für $(0,0)$) und $1 > \theta$
  (für $(1,0)$ und $(0,1)$), also
  \[
    \theta \in [0, 1[, \qquad \text{z.B. } \theta = 0{,}5 .
  \]

  \item XOR verlangt gleichzeitig
  \[
    0 \leq \theta, \qquad w_1 > \theta, \qquad w_2 > \theta,
    \qquad w_1 + w_2 \leq \theta .
  \]
  Aus den mittleren Bedingungen folgt
  \[
    w_1 + w_2 > 2\theta \geq \theta ,
  \]
  im Widerspruch zu $w_1 + w_2 \leq \theta$. XOR ist also
  \textbf{nicht linear separierbar}.
\end{enumerate}
}

% -------------------------------------------------------------------
% 4) XOR mit Hidden-Layer
% -------------------------------------------------------------------
\aufgabe{
Gegeben sei ein Netz mit zwei Hidden-Neuronen $a$ und $b$ sowie einem
Ausgabeneuron $c$. Alle Neuronen verwenden die Heaviside-Funktion, die
Bias-Gewichte sind in den Formeln bereits enthalten:
\[
  o_a = \text{heaviside}(x_1 + x_2 - 0{,}5), \qquad
  o_b = \text{heaviside}(x_1 + x_2 - 1{,}5),
\]
\[
  y = \text{heaviside}(o_a - 2\,o_b - 0{,}5)
\]
Prüfen Sie durch Rechnung für alle vier Eingaben, ob das Netz die
XOR-Funktion realisiert.
}

\loesung{
\[
  \begin{array}{c|c|c|c|c|c}
    x_1 & x_2 & o_a & o_b & o_a - 2o_b - 0{,}5 & y\\\hline
    0 & 0 & 0 & 0 & -0{,}5 & 0\\
    0 & 1 & 1 & 0 & \phantom{-}0{,}5 & 1\\
    1 & 0 & 1 & 0 & \phantom{-}0{,}5 & 1\\
    1 & 1 & 1 & 1 & -1{,}5 & 0
  \end{array}
\]
Beispielrechnung für $\mathbf{x} = (1,1)$:
\[
  o_a = \text{heaviside}(1 + 1 - 0{,}5) = 1, \qquad
  o_b = \text{heaviside}(1 + 1 - 1{,}5) = 1
\]
\[
  y = \text{heaviside}(1 - 2 - 0{,}5) = \text{heaviside}(-1{,}5) = 0
\]
Die Ausgabespalte entspricht genau der XOR-Funktion. Neuron $a$ realisiert
\textbf{OR}, Neuron $b$ realisiert \textbf{AND}, das Ausgabeneuron bildet
$\text{OR} \wedge \neg\text{AND}$.
}

% -------------------------------------------------------------------
% 5) Aktivierungsfunktionen und Ableitungen
% -------------------------------------------------------------------
\aufgabe{
Gegeben sind die Aktivierungsfunktionen
\[
  \text{sig}(ax) = \frac{1}{1+\exp(-ax)}, \qquad
  \tanh(ax) = \frac{e^{ax}-e^{-ax}}{e^{ax}+e^{-ax}} .
\]

\begin{enumerate}[label=\alph*)]
  \item Berechnen Sie $\text{sig}(0)$, $\text{sig}'(0)$, $\tanh(0)$ und
        $\tanh'(0)$.
  \item Für ein Neuron wurde bereits $\text{sig}(\mathrm{net}) = 0{,}8$
        berechnet. Bestimmen Sie $\text{sig}'(\mathrm{net})$ ohne erneutes
        Auswerten der Exponentialfunktion.
  \item Berechnen Sie die Ausgabe eines Neurons mit $a = 5$, $x = 0{,}2$
        und Sigmoid-Aktivierung sowie die zugehörige Ableitung.
        Hinweis: $\exp(-1) \approx 0{,}368$.
\end{enumerate}
}

\loesung{
\begin{enumerate}[label=\alph*)]
  \item
  \[
    \text{sig}(0) = \frac{1}{1+1} = 0{,}5, \qquad
    \text{sig}'(0) = 0{,}5\,(1-0{,}5) = 0{,}25
  \]
  \[
    \tanh(0) = 0, \qquad
    \tanh'(0) = (1+0)(1-0) = 1
  \]

  \item Mit $\text{sig}'(t) = \text{sig}(t)\bigl(1-\text{sig}(t)\bigr)$:
  \[
    \text{sig}'(\mathrm{net}) = 0{,}8\cdot(1-0{,}8) = 0{,}8\cdot 0{,}2
    = \textbf{0{,}16}
  \]

  \item Mit $ax = 5\cdot 0{,}2 = 1$:
  \[
    \text{sig}(1) = \frac{1}{1+0{,}368} \approx 0{,}731
  \]
  \[
    \text{sig}'(1) \approx 0{,}731\cdot(1-0{,}731) = 0{,}731\cdot 0{,}269
    \approx 0{,}197
  \]
\end{enumerate}
}

% -------------------------------------------------------------------
% 6) MLP: Vorwaertsberechnung in Matrixform
% -------------------------------------------------------------------
\aufgabe{
Ein MLP mit einem Hidden-Layer (Sigmoid) und linearer Ausgabeschicht
besitzt die Gewichtsmatrizen
\[
  W^{(1)} =
  \begin{pmatrix}
    1 & 0\\
    -1 & 1\\
    0{,}5 & -0{,}5
  \end{pmatrix}
  \in \mathbb{R}^{3\times 2},
  \qquad
  W^{(2)} = \begin{pmatrix} 1 & -1 & 2 \end{pmatrix} \in \mathbb{R}^{1\times 3}.
\]
Berechnen Sie die Netzausgabe
$\hat{y} = W^{(2)}\,\text{sig}\bigl(W^{(1)}\mathbf{x}\bigr)$
für $\mathbf{x} = \begin{pmatrix} 1 \\ 2 \end{pmatrix}$.
Hinweis: $\text{sig}(1) \approx 0{,}731$, $\text{sig}(-0{,}5) \approx 0{,}378$.
}

\loesung{
Erste Schicht:
\[
  \mathbf{b} = W^{(1)}\mathbf{x} =
  \begin{pmatrix}
    1\cdot 1 + 0\cdot 2\\
    -1\cdot 1 + 1\cdot 2\\
    0{,}5\cdot 1 - 0{,}5\cdot 2
  \end{pmatrix}
  =
  \begin{pmatrix} 1 \\ 1 \\ -0{,}5 \end{pmatrix}
\]
Komponentenweise Sigmoid-Aktivierung:
\[
  O^{(1)} = \text{sig}(\mathbf{b}) \approx
  \begin{pmatrix} 0{,}731 \\ 0{,}731 \\ 0{,}378 \end{pmatrix}
\]
Lineare Ausgabeschicht:
\[
  \hat{y} = 1\cdot 0{,}731 + (-1)\cdot 0{,}731 + 2\cdot 0{,}378
  = \textbf{0{,}756}
\]
}

% -------------------------------------------------------------------
% 7) Lineare Aktivierung: Kollaps der Schichten
% -------------------------------------------------------------------
\aufgabe{
Ein Netz mit zwei Schichten verwendet ausschließlich \textbf{lineare}
Aktivierung, es gilt also $\hat{\mathbf{y}} = W^{(2)}W^{(1)}\mathbf{x}$ mit
\[
  W^{(1)} = \begin{pmatrix} 1 & 2\\ 0 & 1 \end{pmatrix},
  \qquad
  W^{(2)} = \begin{pmatrix} 1 & -1\\ 2 & 0 \end{pmatrix}.
\]

\begin{enumerate}[label=\alph*)]
  \item Bestimmen Sie die Matrix $W = W^{(2)}W^{(1)}$ des äquivalenten
        einlagigen Netzes.
  \item Berechnen Sie $\hat{\mathbf{y}}$ für
        $\mathbf{x} = \begin{pmatrix} 1 \\ 1 \end{pmatrix}$ einmal
        schichtweise und einmal über $W$.
  \item Welche Konsequenz hat das Ergebnis für die Wahl der
        Aktivierungsfunktion?
\end{enumerate}
}

\loesung{
\begin{enumerate}[label=\alph*)]
  \item
  \[
    W = W^{(2)}W^{(1)} =
    \begin{pmatrix} 1 & -1\\ 2 & 0 \end{pmatrix}
    \begin{pmatrix} 1 & 2\\ 0 & 1 \end{pmatrix}
    =
    \begin{pmatrix} 1 & 1\\ 2 & 4 \end{pmatrix}
  \]

  \item Schichtweise:
  \[
    W^{(1)}\mathbf{x} = \begin{pmatrix} 3 \\ 1 \end{pmatrix},
    \qquad
    W^{(2)}\begin{pmatrix} 3 \\ 1 \end{pmatrix}
    = \begin{pmatrix} 2 \\ 6 \end{pmatrix}
  \]
  Über $W$:
  \[
    W\mathbf{x} = \begin{pmatrix} 1+1 \\ 2+4 \end{pmatrix}
    = \begin{pmatrix} 2 \\ 6 \end{pmatrix}
  \]

  \item Mehrere lineare Schichten lassen sich stets zu \textbf{einer}
  Matrix zusammenfassen. Zusätzliche Schichten bringen ohne
  \textbf{nichtlineare} Aktivierung keinen Zugewinn an Modellkomplexität.
\end{enumerate}
}

% -------------------------------------------------------------------
% 8) Gradientenabstieg
% -------------------------------------------------------------------
\aufgabe{
\begin{enumerate}[label=\alph*)]
  \item Gegeben sei $L(w) = (w-3)^2$ mit Startwert $w_0 = 0$ und Lernrate
        $\alpha = 0{,}2$. Führen Sie drei Schritte des Gradientenabstiegs
        $w_{i+1} = w_i - \alpha\,L'(w_i)$ durch.
  \item Zeigen Sie, dass $w_{i+1} - 3 = (1-2\alpha)(w_i - 3)$ gilt, und
        bestimmen Sie, für welche $\alpha$ das Verfahren nicht mehr
        konvergiert.
  \item Führen Sie für $L(w_1,w_2) = w_1^2 + 4w_2^2$ mit
        $\mathbf{w}_0 = (1,1)^{\mathsf{T}}$ und $\alpha = 0{,}1$ zwei
        Schritte durch.
\end{enumerate}
}

\loesung{
\begin{enumerate}[label=\alph*)]
  \item Es gilt $L'(w) = 2(w-3)$:
  \[
    w_1 = 0 - 0{,}2\cdot(-6) = 1{,}2
  \]
  \[
    w_2 = 1{,}2 - 0{,}2\cdot(-3{,}6) = 1{,}92
  \]
  \[
    w_3 = 1{,}92 - 0{,}2\cdot(-2{,}16) = 2{,}352
  \]
  Die Folge nähert sich dem Minimum $w = 3$.

  \item Einsetzen liefert
  \[
    w_{i+1} - 3 = w_i - 2\alpha(w_i-3) - 3 = (1-2\alpha)(w_i-3).
  \]
  Konvergenz erfordert $|1-2\alpha| < 1$, also $0 < \alpha < 1$. Für
  $\alpha \geq 1$ divergiert das Verfahren bzw. springt zwischen zwei
  Werten hin und her.

  \item Gradient: $\nabla L = (2w_1,\; 8w_2)^{\mathsf{T}}$.
  \[
    \mathbf{w}_1 = (1,1)^{\mathsf{T}} - 0{,}1\cdot(2,8)^{\mathsf{T}}
    = (0{,}8;\; 0{,}2)^{\mathsf{T}}
  \]
  \[
    \nabla L(\mathbf{w}_1) = (1{,}6;\; 1{,}6)^{\mathsf{T}}
    \;\Rightarrow\;
    \mathbf{w}_2 = (0{,}64;\; 0{,}04)^{\mathsf{T}}
  \]
  Die steilere Richtung $w_2$ wird deutlich schneller verkleinert.
\end{enumerate}
}

% -------------------------------------------------------------------
% 9) Loss, Batch- und Stochastic Gradient Descent
% -------------------------------------------------------------------
\aufgabe{
Ein Modell mit einem einzigen Gewicht berechnet $\hat{y} = w\cdot x$. Die
Loss Function lautet
\[
  L = \sum_{(x_i,y_i)\in D}(\hat{y}_i - y_i)^2 .
\]
Gegeben sind $D = \{(1,2),\,(2,3),\,(3,7)\}$, $w = 1$ und $\alpha = 0{,}01$.

\begin{enumerate}[label=\alph*)]
  \item Berechnen Sie $L$.
  \item Berechnen Sie $\dfrac{\partial L}{\partial w}$ und führen Sie einen
        Schritt des Gradientenabstiegs über alle Beispiele durch.
  \item Führen Sie stattdessen einen SGD-Schritt nur mit dem Beispiel
        $(3,7)$ durch und vergleichen Sie.
\end{enumerate}
}

\loesung{
\begin{enumerate}[label=\alph*)]
  \item Vorhersagen: $\hat{y} = 1,\,2,\,3$. Damit
  \[
    L = (1-2)^2 + (2-3)^2 + (3-7)^2 = 1 + 1 + 16 = 18 .
  \]

  \item Ableitung:
  \[
    \frac{\partial L}{\partial w}
    = \sum_i 2(\hat{y}_i-y_i)\,x_i
    = 2\bigl((-1)\cdot 1 + (-1)\cdot 2 + (-4)\cdot 3\bigr)
    = -30
  \]
  \[
    w_{\text{neu}} = 1 - 0{,}01\cdot(-30) = \textbf{1{,}3}
  \]

  \item Nur mit $(3,7)$:
  \[
    \frac{\partial L}{\partial w} = 2(3-7)\cdot 3 = -24
    \;\Rightarrow\;
    w_{\text{neu}} = 1 + 0{,}24 = \textbf{1{,}24}
  \]
  Der SGD-Schritt ist deutlich billiger zu berechnen, approximiert die
  Richtung des vollen Gradienten aber nur.
\end{enumerate}
}

% -------------------------------------------------------------------
% 10) Backpropagation: Kette mit Sigmoid
% -------------------------------------------------------------------
\aufgabe{
Gegeben sei ein Netz mit zwei Eingängen, einem versteckten Neuron $h$
(Sigmoid) und einem Ausgabeneuron $o$ (lineare Aktivierung, also
$\varphi(x)=x$ und $\varphi'(x)=1$):
\[
  w_{h,1} = 0{,}5, \qquad w_{h,2} = -0{,}5, \qquad w_{o,h} = 1
\]
Trainingsbeispiel: $x_1 = 1$, $x_2 = -1$, $y = 0{,}5$, Lernrate
$\alpha = 0{,}5$. Verwendet werden
\[
  \delta_j =
  \begin{cases}
    \varphi'(\mathrm{net}_j)(o_j-y) & j \text{ Ausgabeneuron}\\[0.2cm]
    \varphi'(\mathrm{net}_j)\sum_k \delta_k w_{k,j} & \text{sonst}
  \end{cases}
  \qquad
  \Delta w_{i,j} = -\alpha\,\delta_i\,o_j
\]

\begin{enumerate}[label=\alph*)]
  \item Führen Sie die Forward Propagation durch und bestimmen Sie $L$.
  \item Berechnen Sie die Fehlersignale $\delta_o$ und $\delta_h$.
  \item Berechnen Sie alle Gewichtsänderungen und die neuen Gewichte.
\end{enumerate}
Hinweis: $\text{sig}(1) \approx 0{,}731$.
}

\loesung{
\begin{enumerate}[label=\alph*)]
  \item Verstecktes Neuron:
  \[
    \mathrm{net}_h = 0{,}5\cdot 1 + (-0{,}5)\cdot(-1) = 1,
    \qquad
    o_h = \text{sig}(1) \approx 0{,}731
  \]
  Ausgabeneuron:
  \[
    \mathrm{net}_o = 1\cdot 0{,}731 = 0{,}731,
    \qquad
    \hat{y} = o_o = 0{,}731
  \]
  \[
    L = (\hat{y}-y)^2 = (0{,}731-0{,}5)^2 \approx 0{,}053
  \]

  \item Ausgabeneuron (lineare Aktivierung):
  \[
    \delta_o = 1\cdot(0{,}731 - 0{,}5) = 0{,}231
  \]
  Verstecktes Neuron (Sigmoid):
  \[
    \delta_h = o_h(1-o_h)\cdot \delta_o\,w_{o,h}
    = 0{,}731\cdot 0{,}269\cdot 0{,}231\cdot 1 \approx 0{,}045
  \]

  \item Gewichtsänderungen:
  \[
    \Delta w_{o,h} = -0{,}5\cdot 0{,}231\cdot 0{,}731 \approx -0{,}084
    \;\Rightarrow\; w_{o,h} \approx 0{,}916
  \]
  \[
    \Delta w_{h,1} = -0{,}5\cdot 0{,}045\cdot 1 \approx -0{,}023
    \;\Rightarrow\; w_{h,1} \approx 0{,}477
  \]
  \[
    \Delta w_{h,2} = -0{,}5\cdot 0{,}045\cdot(-1) \approx 0{,}023
    \;\Rightarrow\; w_{h,2} \approx -0{,}477
  \]
  Die Ausgabe war zu groß, daher werden die Gewichte so verändert, dass
  $\hat{y}$ sinkt.
\end{enumerate}
}

% -------------------------------------------------------------------
% 11) Backpropagation: zwei Ausgabeneuronen
% -------------------------------------------------------------------
\aufgabe{
Ein Netz besitzt einen Eingang $x$, ein verstecktes Neuron $h$ und
\textbf{zwei} Ausgabeneuronen $o_1$ und $o_2$. Alle Neuronen verwenden die
Identität als Aktivierung, es gilt also $\varphi'(x) = 1$.
\[
  w_{h,x} = 0{,}5, \qquad w_{o_1,h} = 1, \qquad w_{o_2,h} = -2
\]
Trainingsbeispiel: $x = 2$ mit den Sollwerten $y_1 = 0$ und $y_2 = -1$,
Lernrate $\alpha = 0{,}1$.

\begin{enumerate}[label=\alph*)]
  \item Berechnen Sie die Forward Propagation sowie den Gesamtfehler
        $L = (o_1-y_1)^2 + (o_2-y_2)^2$.
  \item Berechnen Sie $\delta_{o_1}$, $\delta_{o_2}$ und daraus $\delta_h$.
        Beachten Sie die Summe über \textbf{alle} nachfolgenden Neuronen.
  \item Berechnen Sie alle Gewichtsänderungen und die neuen Gewichte.
\end{enumerate}
}

\loesung{
\begin{enumerate}[label=\alph*)]
  \item Verstecktes Neuron:
  \[
    \mathrm{net}_h = 0{,}5\cdot 2 = 1 \;\Rightarrow\; o_h = 1
  \]
  Ausgabeneuronen:
  \[
    \mathrm{net}_{o_1} = 1\cdot 1 = 1 \;\Rightarrow\; o_1 = 1,
    \qquad
    \mathrm{net}_{o_2} = -2\cdot 1 = -2 \;\Rightarrow\; o_2 = -2
  \]
  \[
    L = (1-0)^2 + (-2-(-1))^2 = 1 + 1 = 2
  \]

  \item Ausgabeschicht:
  \[
    \delta_{o_1} = 1\cdot(1-0) = 1,
    \qquad
    \delta_{o_2} = 1\cdot(-2+1) = -1
  \]
  Verstecktes Neuron:
  \[
    \delta_h = 1\cdot\bigl(\delta_{o_1}w_{o_1,h} + \delta_{o_2}w_{o_2,h}\bigr)
    = 1\cdot 1 + (-1)\cdot(-2) = 3
  \]

  \item Mit $\Delta w_{i,j} = -\alpha\,\delta_i\,o_j$:
  \[
    \Delta w_{o_1,h} = -0{,}1\cdot 1\cdot 1 = -0{,}1
    \;\Rightarrow\; w_{o_1,h} = 0{,}9
  \]
  \[
    \Delta w_{o_2,h} = -0{,}1\cdot(-1)\cdot 1 = 0{,}1
    \;\Rightarrow\; w_{o_2,h} = -1{,}9
  \]
  \[
    \Delta w_{h,x} = -0{,}1\cdot 3\cdot 2 = -0{,}6
    \;\Rightarrow\; w_{h,x} = -0{,}1
  \]
  Das versteckte Neuron sammelt die Fehlersignale \textbf{beider}
  Ausgabeneuronen ein.
\end{enumerate}
}

% -------------------------------------------------------------------
% 12) Backpropagation: vollstaendiges Netz mit Bias
% -------------------------------------------------------------------
\aufgabe{
Ein Netz besitzt zwei Eingänge, zwei versteckte Neuronen $h_1, h_2$ und ein
Ausgabeneuron $o$. Alle Neuronen verwenden die \textbf{Sigmoid}-Funktion,
jede Schicht besitzt ein Bias-Neuron mit konstantem Wert $1$. Die Gewichte
lauten
\[
  \begin{array}{lll}
    w_{1,1} = 0{,}5, & w_{1,2} = 0{,}5, & w_{1,b} = 0\\
    w_{2,1} = -0{,}5, & w_{2,2} = -0{,}5, & w_{2,b} = 0\\
    w_{o,1} = 1, & w_{o,2} = -1, & w_{o,b} = 0
  \end{array}
\]
Trainingsbeispiel: $x_1 = 1$, $x_2 = -1$, $y = 1$, Lernrate
$\alpha = 0{,}4$.

\begin{enumerate}[label=\alph*)]
  \item Berechnen Sie die Ausgabe $\hat{y}$ und den Fehler $L$.
  \item Berechnen Sie $\delta_o$, $\delta_1$ und $\delta_2$.
  \item Berechnen Sie die Gewichtsänderungen der Ausgabeschicht sowie die
        von $h_1$.
\end{enumerate}
Hinweis: $\text{sig}(0) = 0{,}5$ und $\text{sig}'(0) = 0{,}25$.
}

\loesung{
\begin{enumerate}[label=\alph*)]
  \item Hidden-Layer:
  \[
    \mathrm{net}_1 = 0{,}5\cdot 1 + 0{,}5\cdot(-1) + 0 = 0
    \;\Rightarrow\; o_1 = 0{,}5
  \]
  \[
    \mathrm{net}_2 = -0{,}5\cdot 1 - 0{,}5\cdot(-1) + 0 = 0
    \;\Rightarrow\; o_2 = 0{,}5
  \]
  Ausgabeneuron:
  \[
    \mathrm{net}_o = 1\cdot 0{,}5 + (-1)\cdot 0{,}5 + 0 = 0
    \;\Rightarrow\; \hat{y} = 0{,}5
  \]
  \[
    L = (0{,}5-1)^2 = 0{,}25
  \]

  \item Ausgabeneuron:
  \[
    \delta_o = \hat{y}(1-\hat{y})(\hat{y}-y)
    = 0{,}25\cdot(0{,}5-1) = -0{,}125
  \]
  Versteckte Neuronen:
  \[
    \delta_1 = o_1(1-o_1)\,\delta_o\,w_{o,1}
    = 0{,}25\cdot(-0{,}125)\cdot 1 = -0{,}03125
  \]
  \[
    \delta_2 = o_2(1-o_2)\,\delta_o\,w_{o,2}
    = 0{,}25\cdot(-0{,}125)\cdot(-1) = 0{,}03125
  \]

  \item Ausgabeschicht mit $\Delta w_{i,j} = -\alpha\,\delta_i\,o_j$:
  \[
    \Delta w_{o,1} = -0{,}4\cdot(-0{,}125)\cdot 0{,}5 = 0{,}025
    \;\Rightarrow\; w_{o,1} = 1{,}025
  \]
  \[
    \Delta w_{o,2} = -0{,}4\cdot(-0{,}125)\cdot 0{,}5 = 0{,}025
    \;\Rightarrow\; w_{o,2} = -0{,}975
  \]
  \[
    \Delta w_{o,b} = -0{,}4\cdot(-0{,}125)\cdot 1 = 0{,}05
    \;\Rightarrow\; w_{o,b} = 0{,}05
  \]
  Neuron $h_1$:
  \[
    \Delta w_{1,1} = -0{,}4\cdot(-0{,}03125)\cdot 1 = 0{,}0125
    \;\Rightarrow\; w_{1,1} = 0{,}5125
  \]
  \[
    \Delta w_{1,2} = -0{,}4\cdot(-0{,}03125)\cdot(-1) = -0{,}0125
    \;\Rightarrow\; w_{1,2} = 0{,}4875
  \]
  \[
    \Delta w_{1,b} = -0{,}4\cdot(-0{,}03125)\cdot 1 = 0{,}0125
    \;\Rightarrow\; w_{1,b} = 0{,}0125
  \]
  Die Ausgabe war zu klein, alle Änderungen wirken in Richtung eines
  größeren $\hat{y}$.
\end{enumerate}
}


% -------------------------------------------------------------------
% 13) Backpropagation
% -------------------------------------------------------------------

\aufgabe{
Betrachtet wird ein neuronales Netz mit zwei Eingabeneuronen, einem
versteckten Layer aus zwei Neuronen mit ReLU-Aktivierung
$\varphi : \mathbb{R} \to \mathbb{R},\ x \mapsto \max(0,x)$ und einem
Ausgabeneuron mit Identität
$\varphi : \mathbb{R} \to \mathbb{R},\ x \mapsto x$. Jede Schicht besitzt
ein \textbf{Bias-Neuron} mit konstantem Wert $1$. Alle Gewichte der
Eingabe-- und Ausgabeverbindungen sind positiv.

\begin{center}
\begin{tikzpicture}[scale=1,
    neuron/.style={circle, draw, minimum size=0.8cm},
    bias/.style={circle, draw, fill=black!8, minimum size=0.7cm},
    >=stealth
]
  \node[neuron] (x1) at (0,1.4)  {$x_1$};
  \node[neuron] (x2) at (0,-0.6) {$x_2$};
  \node[bias]   (b1) at (0,-2.4) {$1$};
  \node[neuron] (h1) at (3,1.4)  {$h_1$};
  \node[neuron] (h2) at (3,-0.6) {$h_2$};
  \node[bias]   (b2) at (3,-2.4) {$1$};
  \node[neuron] (y)  at (6,0.4)  {$\hat{y}$};
  \draw[->] (x1) -- node[pos=0.8, above] {$0{,}5$} (h1);
  \draw[->] (x1) -- node[pos=0.8, right=1pt] {$0{,}3$} (h2);
  \draw[->] (x2) -- node[pos=0.85, left=2pt] {$0{,}4$} (h1);
  \draw[->] (x2) -- node[pos=0.8, below] {$0{,}7$} (h2);
  \draw[->] (b1) -- node[pos=0.85, right] {$0{,}2$} (h1);
  \draw[->] (b1) -- node[pos=0.8, below] {$-0{,}1$} (h2);
  \draw[->] (h1) -- node[above] {$0{,}8$} (y);
  \draw[->] (h2) -- node[below] {$0{,}5$} (y);
  \draw[->] (b2) -- node[below] {$0{,}3$} (y);
  \node[draw=none, above] at (0,1.9) {Eingabe};
  \node[draw=none, above] at (3,1.9) {Versteckt};
  \node[draw=none, above] at (6,1.9) {Ausgabe};
\end{tikzpicture}
\end{center}

Gegeben sei der Input $\mathbf{x} = (1,2)^{\mathsf{T}}$ mit Sollwert $y = 2$
und Lernrate $\alpha = 0{,}1$. Das Bias-Gewicht eines Neurons wird
behandelt wie ein gewöhnliches Gewicht mit konstanter Eingabe $o = 1$.

\begin{enumerate}[label=\alph*)]
  \item Führen Sie die Vorwärtsrechnung durch und bestimmen Sie
        $h_1$, $h_2$ und $\hat{y}$.
  \item Berechnen Sie die Fehlersignale $\delta_j$ nach
        \[
          \delta_j =
          \begin{cases}
            \varphi'(\mathrm{net}_j)\,(\hat{y} - y),
              & j \text{ Ausgabeneuron},\\[0.2cm]
            \varphi'(\mathrm{net}_j)\displaystyle\sum_k \delta_k\,w_{jk},
              & \text{sonst}.
          \end{cases}
        \]
  \item Berechnen Sie die neuen Gewichte \textbf{einschließlich der
        Bias-Gewichte} nach einer Iteration des Gradientenabstiegs
        \[
          w_{ij}' = w_{ij} + \Delta w_{ij},
          \qquad
          \Delta w_{ij} = -\alpha\,\delta_j\,o_i .
        \]
\end{enumerate}
}

\loesung{
\begin{enumerate}[label=\alph*)]
  \item \textbf{Vorwärtsrechnung.} Das Bias-Neuron liefert den konstanten
  Summanden $1\cdot b_j$. Da beide Netzeingaben positiv sind, gilt
  $\varphi(\mathrm{net}) = \mathrm{net}$ für die ReLU-Neuronen:
  \begin{align*}
    h_1 &= \varphi(1\cdot 0{,}5 + 2\cdot 0{,}4 + 1\cdot 0{,}2)
         = \varphi(1{,}5) = 1{,}5,\\
    h_2 &= \varphi(1\cdot 0{,}3 + 2\cdot 0{,}7 + 1\cdot(-0{,}1))
         = \varphi(1{,}6) = 1{,}6,\\
    \hat{y} &= 0{,}8\cdot 1{,}5 + 0{,}5\cdot 1{,}6 + 1\cdot 0{,}3 = 2{,}3 .
  \end{align*}

  \item \textbf{Fehlersignale.} Wegen $\varphi'(\mathrm{net}_j) = 1$ für
  die Identität und für ReLU bei positiver Netzeingabe:
  \begin{align*}
    \delta_{\hat{y}} &= 1\cdot(2{,}3 - 2) = 0{,}3,\\
    \delta_{h_1} &= 1\cdot \delta_{\hat{y}}\,w_{h_1\hat{y}}
                  = 0{,}3\cdot 0{,}8 = 0{,}24,\\
    \delta_{h_2} &= 1\cdot \delta_{\hat{y}}\,w_{h_2\hat{y}}
                  = 0{,}3\cdot 0{,}5 = 0{,}15 .
  \end{align*}
  Die Fehlersignale der Bias-Neuronen müssen nicht berechnet werden, da
  sie keine eingehenden Verbindungen besitzen.

  \item \textbf{Gewichtsaktualisierung.} Für die Bias-Gewichte ist die
  zugehörige Eingabe $o = 1$.

  \textbf{Versteckt $\to$ Ausgabe (inkl. Bias):}
  \begin{align*}
    \Delta w_{h_1\hat{y}} &= -0{,}1\cdot 0{,}3\cdot 1{,}5 = -0{,}045,
      & w_{h_1\hat{y}}' &= 0{,}8 - 0{,}045 = 0{,}755,\\
    \Delta w_{h_2\hat{y}} &= -0{,}1\cdot 0{,}3\cdot 1{,}6 = -0{,}048,
      & w_{h_2\hat{y}}' &= 0{,}5 - 0{,}048 = 0{,}452,\\
    \Delta b_{\hat{y}} &= -0{,}1\cdot 0{,}3\cdot 1 = -0{,}03,
      & b_{\hat{y}}' &= 0{,}3 - 0{,}03 = 0{,}27 .
  \end{align*}

  \textbf{Eingabe $\to$ Versteckt (inkl. Bias):}
  \begin{align*}
    \Delta w_{x_1 h_1} &= -0{,}1\cdot 0{,}24\cdot 1 = -0{,}024,
      & w_{x_1 h_1}' &= 0{,}5 - 0{,}024 = 0{,}476,\\
    \Delta w_{x_2 h_1} &= -0{,}1\cdot 0{,}24\cdot 2 = -0{,}048,
      & w_{x_2 h_1}' &= 0{,}4 - 0{,}048 = 0{,}352,\\
    \Delta b_{h_1} &= -0{,}1\cdot 0{,}24\cdot 1 = -0{,}024,
      & b_{h_1}' &= 0{,}2 - 0{,}024 = 0{,}176,\\
    \Delta w_{x_1 h_2} &= -0{,}1\cdot 0{,}15\cdot 1 = -0{,}015,
      & w_{x_1 h_2}' &= 0{,}3 - 0{,}015 = 0{,}285,\\
    \Delta w_{x_2 h_2} &= -0{,}1\cdot 0{,}15\cdot 2 = -0{,}03,
      & w_{x_2 h_2}' &= 0{,}7 - 0{,}03 = 0{,}67,\\
    \Delta b_{h_2} &= -0{,}1\cdot 0{,}15\cdot 1 = -0{,}015,
      & b_{h_2}' &= -0{,}1 - 0{,}015 = -0{,}115 .
  \end{align*}

  Da $\hat{y} = 2{,}3 > 2 = y$ war, ist der Fehler positiv und sämtliche
  Gewichte samt Bias werden verkleinert, sodass $\hat{y}$ sinkt.
\end{enumerate}
}



% -------------------------------------------------------------------
% 14) Netzgroesse, Batches und Epochen
% -------------------------------------------------------------------
\aufgabe{
Ein MLP besitzt $4$ Eingangsneuronen, zwei Hidden-Layer mit $6$ bzw. $5$
Neuronen und $3$ Ausgabeneuronen. Jede Schicht außer der Ausgabeschicht
besitzt zusätzlich ein Bias-Neuron.

\begin{enumerate}[label=\alph*)]
  \item Geben Sie die Dimensionen von $W^{(1)}$, $W^{(2)}$ und $W^{(3)}$ an
        und berechnen Sie die Gesamtzahl der Gewichte.
  \item Das Netz wird mit $50\,000$ Beispielen bei einer Batch-Größe von
        $200$ über $30$ Epochen trainiert. Wie viele Iterationsschritte
        (Gewichts-Updates) ergeben sich?
\end{enumerate}
}

\loesung{
\begin{enumerate}[label=\alph*)]
  \item Jede Matrix bildet die um das Bias-Neuron erweiterte Vorgängerschicht
  auf die Folgeschicht ab:
  \[
    W^{(1)} \in \mathbb{R}^{6\times 5}, \qquad
    W^{(2)} \in \mathbb{R}^{5\times 7}, \qquad
    W^{(3)} \in \mathbb{R}^{3\times 6}
  \]
  \[
    6\cdot 5 + 5\cdot 7 + 3\cdot 6 = 30 + 35 + 18 = \textbf{83}
  \]

  \item Iterationen pro Epoche:
  \[
    \frac{50\,000}{200} = 250
  \]
  Insgesamt:
  \[
    250\cdot 30 = \textbf{7500} \text{ Gewichts-Updates}
  \]
\end{enumerate}
}


\end{document}